% !TeX root = ../../Book5.tex
% 纲要标题：## 第4章　正规化、离散赋值与曲线奇点
% 纲要位置：几何学计划纲要.md，第 5147 行。
% 正文状态：初稿；尚未正式审核。

% 纲要内容开始
% 
% > 本章目标与证明范围：把正规化对象、有限性证明和局部曲线计算独立成章。任何一般有限性结论都注明有限型或 Nagata 类条件，不能推广到任意 Noether 环。
% 
% 纲要内容结束

\chapter{正规化、离散赋值与曲线奇点}
\label{g5:ch:04}
\BookChapterNumber{g5:ch:04}

\begin{chapterabstract}
	本章在函数域中补入整元，构造正规化，并分析它怎样改变曲线的局部函数。首先证明一维 Noether 正规局部整环恰为离散赋值环，使消失阶、统一化元与理想分解获得统一的代数描述。随后证明任意域上有限型约化代数的正规化有限：可分部分用迹配对控制分母，不可分部分用有限个系数的开根处理。导子理想定位正规化发生变化的位置，正规化商的长度定义曲线点的 \(\delta\) 不变量。尖点、结点、扩域及非约化模型进一步区分有限性、同胚、同构与正规性。
\end{chapterabstract}

\begin{learninggoals}
	\begin{itemize}
		\item 证明 DVR 的等价刻画，计算统一化元、函数的阶、分式理想及局部商的长度。
		\item 从各点的 DVR 结构恢复 Dedekind 整环的非零理想分解，区分局部主性与整体主性。
		\item 在分式域或总分式环中构造正规化，证明局部化相容及任意基域上有限型代数的正规化有限。
		\item 运用迹配对与有限系数开根，说明可分和不可分两种有限性论证的作用与条件。
		\item 计算导子、正规化纤维与 \(\delta\) 不变量，辨认尖点、结点以及基域扩张和幂零结构带来的差异。
	\end{itemize}
\end{learninggoals}

在参数直线上，函数 \(t\) 本身正则；把它映到尖点 \((x,y)=(t^2,t^3)\) 后，原点附近的函数环却不包含 \(t\)。这个缺口不能用增加一个点来弥补，因为参数化已经是点集上的双射。需要改变的是允许作为正则函数的元素：把分式 \(y/x\) 所代表的 \(t\) 补入函数环。

\ProseParagraphBreak

结点呈现另一种情况。设基域特征不为二，参数化
\[
	t\longmapsto\bigl(t^2-1,t(t^2-1)\bigr)
\]
把 \(t=1\) 与 \(t=-1\) 送到同一个点。原曲线的函数必须在这两个参数值取相同的值，补入 \(t\) 后，这个相等条件被解除，两个原像可以分开。尖点缺少的一次项与结点缺少的取值差，都能用坐标环包含 \(A\subseteq A^\nu\) 及商模 \(A^\nu/A\) 表达；导子则把发生变化的位置记录为原环中的理想。

\ProseParagraphBreak

第三章已建立整扩张、Noether 正规化、维数与有限长度。本章据此完成普通正规化的构造和有限性证明，再转入具体曲线计算。局部结构以 DVR 为核心，不把统一化元当作整个函数域的生成元，也不预先使用完备化或光滑态射理论。关于一般 Noether 环的有限性边界，最后一节会明确指出需要补充的条件；本章的主要有限性定理则在任意基域上的有限型范围内完整证明。

% !TeX root = ../../Book5.tex
% 纲要标题：### 4.1 一维正规局部结构
% 纲要位置：几何学计划纲要.md，第 5151 行。
% 正文状态：初稿；尚未正式审核。

% 纲要内容开始
% 
% * 离散赋值环的等价刻画
% * 统一化元、阶和分式理想
% * Dedekind 整环的局部观点
% * 正规曲线的局部参数
% 
% 纲要内容结束

\section{一维正规局部结构}
\label{g5:sec:0401}

在直线原点附近，非零有理函数可以唯一写成 \(t^n u\)，其中 \(u\) 在原点既无零点也无极点。整数 \(n\) 记录全部消失或极点信息，剩下的 \(u\) 是局部环中的单位。这种表达不只属于直线：一个一维 Noether 局部整环只要整闭，就有完全相同的结构。因而在正规曲线的一点附近，理想、函数的消失和有限长度都可以化为整数的计算。

\ProseParagraphBreak

本节从局部环建立这种阶函数，再把不同点的阶组织为 Dedekind 整环中的理想分解。所用的主要工具是第二章的行列式技巧、Nakayama 引理和局部检测，以及第三章的整闭性与零维 Noether 环。DVR 与理想可逆性的基本刻画可对照\cite[第11节，定理11.1--11.3，第78--81页]{Matsumura1989}；Dedekind 整环的相应结论见该节定理11.6。

\subsection{离散赋值环的等价刻画}
\label{g5:subsec:0401:01}

先把“阶”应当满足的运算规则写下来。乘积的阶相加；相加时可能抵消最低阶项，但不会产生比两个加数都更低的阶。因此阶自然取值于带有加法和大小关系的整数集。

\begin{definition}
\label{g5:def:discrete-valuation-ring}
	设 \(K\) 为域。一个满射 \(v:K^\times\rightarrow\mathbb Z\) 若满足
	\[
		v(ab)=v(a)+v(b),\qquad
		v(a+b)\geq\min\{v(a),v(b)\}\quad(a+b\ne0),
	\]
	就称为一个规范化的\term[lisanfuzhi]{离散赋值}[discrete valuation]。规定 \(v(0)=+\infty\)，便可把第二个不等式用于任意 \(a,b\in K\)。环
	\[
		R_v=\{a\in K:v(a)\geq0\}
	\]
	称为这个赋值的\term[lisanfuzhihuan]{离散赋值环}[discrete valuation ring]，简称 DVR。
\end{definition}

这里整数值群确为整个 \(\mathbb Z\)，所以赋值非平凡；域本身不属于本节的 DVR。由乘法公式，\(v(1)=v(-1)=0\)，且 \(v(a^{-1})=-v(a)\)。因此 \(R_v\) 的单位正是阶为零的元素，其余非零元素的阶都为正。集合
\[
	\mathfrak m_v=\{a\in K:v(a)>0\}
\]
是理想，且每个不属于它的 \(R_v\) 元素都可逆，故 \(R_v\) 是局部整环。对任意 \(z\in K^\times\)，\(z,z^{-1}\) 至少有一个属于 \(R_v\)，所以 \(K\) 正是 \(R_v\) 的分式域。

\begin{theorem}
\label{g5:thm:dvr-characterizations}
	设 \((R,\mathfrak m)\) 为非域的交换含幺局部整环，\(K=\operatorname{Frac}(R)\)。下列条件等价：
	\begin{equivalences}
		\item \(R\) 为一维 Noether 整闭整环；
		\item \(R\) 为主理想整环；
		\item \(R\) 为离散赋值环；
		\item \(R\) 为 Noether 环，且 \(\mathfrak m\) 为主理想。
	\end{equivalences}
\end{theorem}
\begin{proof}
	先由第一条件证明第四条件。取 \(0\ne a\in\mathfrak m\)。一维局部整环只有两个素理想 \((0)\) 和 \(\mathfrak m\)，所以 \(\sqrt{aR}=\mathfrak m\)。由于 \(\mathfrak m\) 有限生成，存在 \(n\geq1\) 使 \(\mathfrak m^n\subseteq aR\)：若 \(\mathfrak m=(x_1,\ldots,x_r)\)，为每个 \(x_i\) 选一个落入 \(aR\) 的幂，再用足够高次单项式必含其中某个幂即可。

	取满足上述包含的最小 \(n\)，并选
	\[
		b\in\mathfrak m^{n-1}\setminus aR,\qquad z=b/a\in K.
	\]
	于是 \(z\notin R\)，但 \(z\mathfrak m\subseteq R\)。我们断言 \(z\mathfrak m\) 不包含在 \(\mathfrak m\) 中。否则，乘以 \(z\) 就是有限生成 \(R\) 模 \(\mathfrak m\) 上的线性算子；\cref{g5:lem:determinant-trick}给出首一多项式 \(F(T)\in R[T]\)，使
	\[
		F(z)\mathfrak m=0.
	\]
	\(\mathfrak m\) 含非零元素，而这些等式位于域 \(K\) 中，故 \(F(z)=0\)。这样 \(z\) 整于 \(R\)，整闭性迫使 \(z\in R\)，矛盾。

	因此 \(z\mathfrak m\) 是不包含在极大理想中的 \(R\) 理想，必等于 \(R\)。可取 \(\pi\in\mathfrak m\) 使 \(z\pi=1\)。对每个 \(x\in\mathfrak m\)，都有 \(zx\in R\)，于是
	\[
		x=\pi(zx)\in\pi R.
	\]
	所以 \(\mathfrak m=\pi R\)，得到第四条件。

	现在设第四条件成立，写 \(\mathfrak m=\pi R\)，其中 \(\pi\ne0\) 且不是单位。我们证明每个 \(0\ne a\in R\) 都可写为
	\[
		a=u\pi^n,\qquad u\in R^\times,\quad n\geq0.
	\]
	若 \(a\) 不是单位，它属于 \(\pi R\)，可先除以一个 \(\pi\)，然后对商重复。这个过程必须终止。否则令 \(a_j=a/\pi^j\in R\)，便得到上升链
	\[
		a_0R\subseteq a_1R\subseteq a_2R\subseteq\cdots.
	\]
	Noether 性使某一步相等，于是 \(a_{j+1}=ra_j=r\pi a_{j+1}\)。整环中可约去非零的 \(a_{j+1}\)，得到 \(r\pi=1\)，矛盾。表达式中的指数唯一：若 \(u\pi^n=w\pi^m\) 且 \(n<m\)，约去 \(\pi^n\) 后会把单位 \(u\) 写成极大理想中的元素。

	对 \(0\ne I\subseteq R\)，取其中非零元素所对应指数的最小值 \(n\)。一个达到最小值的元素 \(u\pi^n\) 生成 \(\pi^nR\)，而 \(I\) 的其余元素都属于 \(\pi^nR\)。故
	\[
		I=\pi^nR.
	\]
	零理想也是主理想，这证明第四条件蕴含第二条件。反之，主理想整环的每个理想均有限生成，且 \(\mathfrak m\) 为主理想，所以第二条件蕴含第四条件。

	仍由第四条件所得到的唯一指数表达，定义
	\[
		v(a/b)=n(a)-n(b)\qquad(a,b\in R\setminus\{0\}).
	\]
	若 \(a/b=c/d\)，则 \(ad=bc\)，乘积的唯一指数给出 \(n(a)+n(d)=n(b)+n(c)\)，故定义无关。乘法时指数相加。相加时，把两个分式共同提出较小的 \(\pi\) 次幂，剩余两项都在 \(R\) 中，所以阶不小于较小者。又 \(v(\pi)=1\)，且恰在 \(v(x)\geq0\) 时 \(x\in R\)。这便构造了第三条件中的离散赋值。

	最后设第三条件成立。选 \(v(\pi)=1\)。每个非零 \(a\in R\) 都有 \(a=\pi^{v(a)}u\)，其中 \(v(u)=0\)，即 \(u\) 为单位。取最小阶的同样论证表明，每个非零理想都是 \(\pi^nR\)，故 \(R\) Noether。其非零素理想只能为 \(\pi R\)：若素理想包含 \(\pi^n\)，就必须包含 \(\pi\)。因此 \(\dim R=1\)。

	为证明整闭，若 \(x\in K\) 整于 \(R\) 却有 \(v(x)<0\)，则 \(x^{-1}\in\mathfrak m_v\)。把 \(x\) 的首一方程除以最高次幂，就得到
	\[
		1+a_1x^{-1}+\cdots+a_nx^{-n}=0.
	\]
	右侧除 \(1\) 外的各项都属于 \(\mathfrak m_v\)，这会使 \(1\in\mathfrak m_v\)，矛盾。故 \(v(x)\geq0\)，即 \(x\in R\)。第三条件蕴含第一条件，全部等价关系得证。
\end{proof}

证明中的两个有限性步骤承担不同作用：极大理想有限生成，使行列式技巧能把稳定性变成整性；理想升链终止，使一个非零元素不能无限次除以 \(\pi\)。第三章最后的非 Noether 例子已有主极大理想却包含更长的素链，说明仅有极大理想为主理想还不够。

\begin{example}
\label{g5:exm:basic-dvrs}
	设 \(k\) 为域，\(f(t)\in k[t]\) 为不可约多项式。则
	\[
		k[t]_{(f)}
	\]
	是非域局部主理想整环，因而是 DVR；其阶计算有理函数的分子、分母中 \(f\) 的重数差。特别地，\(f=t-a\) 时得到点 \(a\in k\) 处的消失阶。算术例子是 \(\mathbb Z_{(p)}\)，其中 \(p\) 为素数；对应的阶为有理数分子、分母中 \(p\) 的指数差。这里的剩余域分别为 \(k[t]/(f)\) 和 \(\mathbb F_p\)。
\end{example}

\subsection{统一化元、阶和分式理想}
\label{g5:subsec:0401:02}

DVR 的极大理想用一个元素生成，所有非零函数和理想也随之排列成同一列幂次。选择极大理想的生成元，就是选择一把测量消失阶的尺子；换成另一生成元只会乘上单位。

\begin{definition}
\label{g5:def:uniformizer}
	设 \((R,\mathfrak m)\) 为 DVR。满足 \(\mathfrak m=\pi R\) 的元素 \(\pi\) 称为\term[tongyihuayuan]{统一化元}[uniformizing element]。对 \(0\ne f\in\operatorname{Frac}(R)\)，唯一表达
	\[
		f=u\pi^n,\qquad u\in R^\times,\quad n\in\mathbb Z
	\]
	中的整数 \(n\) 称为 \(f\) 的\term[jubuji]{阶}[order]，记为 \(\operatorname{ord}_R(f)\)，并约定 \(\operatorname{ord}_R(0)=+\infty\)。\symidx{\(\operatorname{ord}_R(f)\)：DVR 中非零分式的阶}
\end{definition}

若 \(\pi'\) 也是统一化元，则 \(\pi'=w\pi\)，其中 \(w\) 为单位，因此 \((\pi')^n\) 与 \(\pi^n\) 只差一个单位。阶与选择无关，正是这个 DVR 唯一的规范化离散赋值。其运算满足
\[
	\operatorname{ord}_R(fg)=\operatorname{ord}_R(f)+\operatorname{ord}_R(g),
	\qquad
	\operatorname{ord}_R(f+g)\geq
	\min\{\operatorname{ord}_R(f),\operatorname{ord}_R(g)\}.
\]
若两项的阶不同，不等式必为等号：例如 \(f=\pi^a u\)、\(g=\pi^b w\) 且 \(a<b\)，则
\[
	f+g=\pi^a(u+\pi^{b-a}w),
\]
括号中第一项为单位、第二项在极大理想中，故其和仍为单位。同阶时则可能抵消，这正是只保留不等式的原因。

\begin{definition}
\label{g5:def:fractional-ideal}
	设 \(A\) 为整环，\(K=\operatorname{Frac}(A)\)。一个非零 \(A\) 子模 \(I\subseteq K\)，若存在 \(0\ne d\in A\) 使 \(dI\subseteq A\)，就称为 \(A\) 的\term[fenshilixiang]{分式理想}[fractional ideal]。两个分式理想的乘积由全部有限和 \(\sum_i a_ib_i\)、\(a_i\in I,b_i\in J\) 组成。定义
	\[
		I^{-1}=\{x\in K:xI\subseteq A\}.
	\]
	若 \(II^{-1}=A\)，就称 \(I\) 为\term[kenilixiang]{可逆理想}[invertible ideal]。
\end{definition}

分式理想允许统一的分母，却仍要求所有元素可由同一个非零 \(d\) 清除分母。若 \(A\) Noether，则 \(dI\) 是有限生成的普通理想，所以 \(I\) 也是有限生成 \(A\) 模。集合 \(I^{-1}\) 仍为分式理想：它包含 \(dA\)，因此非零；若 \(0\ne a/b\in I\)、\(a,b\in A\)，则 \(x\in I^{-1}\) 蕴含 \(ax\in bA\subseteq A\)，所以 \(aI^{-1}\subseteq A\)。同样，分式理想的乘积、和与交都仍是分式理想；交非零是因为任意两个非零主分式理想都有非零公共元素。

\begin{proposition}
\label{g5:prop:dvr-fractional-ideals}
	设 \(R\) 为 DVR，\(\pi\) 为统一化元。每个分式理想都唯一写成 \(\pi^nR\)，其中 \(n\in\mathbb Z\)。特别地，每个分式理想均可逆。对整数 \(a,b\)，各项运算如\cref{g5:tab:dvr-ideal-arithmetic}所示。
\end{proposition}
\begin{proof}
	设 \(I\) 为分式理想，取 \(0\ne d\in R\) 使 \(dI\subseteq R\)。已经证明普通非零理想都是统一化元的非负整数次幂，所以 \(dI=\pi^mR\)。写 \(d=u\pi^r\)、\(u\) 为单位，便有 \(I=\pi^{m-r}R\)。唯一性来自最小阶的唯一性。

	因此包含关系为 \(\pi^aR\subseteq\pi^bR\) 当且仅当 \(a\geq b\)。乘积的指数相加，和取包含两者的较大理想，交取较小理想。最后，\(x\pi^aR\subseteq R\) 当且仅当 \(\operatorname{ord}_R(x)\geq-a\)，故 \((\pi^aR)^{-1}=\pi^{-a}R\)，其乘积为 \(R\)。
\end{proof}

\begin{table}[htbp]
	\centering
	\caption{DVR 中非零分式理想的运算}
	\label{g5:tab:dvr-ideal-arithmetic}
	\begin{tabular}{lll}
		\toprule
		运算 & 结果 & 指数的变化 \\
		\midrule
		\(\pi^aR\cdot\pi^bR\) & \(\pi^{a+b}R\) & 相加 \\
		\(\pi^aR+\pi^bR\) & \(\pi^{\min(a,b)}R\) & 取较小者 \\
		\(\pi^aR\cap\pi^bR\) & \(\pi^{\max(a,b)}R\) & 取较大者 \\
		\((\pi^aR)^{-1}\) & \(\pi^{-a}R\) & 变号 \\
		\bottomrule
	\end{tabular}
\end{table}

例如在 \(R=k[t]_{(t)}\) 中，分式理想 \(t^{-2}R\) 包含 \(R\)，其逆为 \(t^2R\)；而 \(t^2R+t^5R=t^2R\)、\(t^2R\cap t^5R=t^5R\)。较高的消失阶意味着对元素施加更强的限制，因此对应较小的理想。

\begin{proposition}
\label{g5:prop:dvr-length}
	设 \((R,\mathfrak m)\) 为 DVR，剩余域为 \(\kappa=R/\mathfrak m\)，\(\pi\) 为统一化元。对 \(n\geq1\)，模 \(R/\pi^nR\) 的长度为 \(n\)。因而对每个 \(0\ne f\in R\)，有
	\[
		\ell_R(R/fR)=\operatorname{ord}_R(f).
	\]
	此外，\(\mathfrak m/\mathfrak m^2\) 是一维 \(\kappa\) 向量空间，\(t\in\mathfrak m\) 是统一化元，当且仅当其在 \(\mathfrak m/\mathfrak m^2\) 中的像非零。
\end{proposition}
\begin{proof}
	考虑滤过
	\[
		R\supset\pi R\supset\cdots\supset\pi^nR.
	\]
	对 \(0\leq j<n\)，乘以 \(\pi^j\) 给出同构
	\[
		R/\pi R\longrightarrow\pi^jR/\pi^{j+1}R.
	\]
	满性来自右侧的表示；若 \(a\pi^j\in\pi^{j+1}R\)，约去非零 \(\pi^j\) 得 \(a\in\pi R\)，证明单射。因此滤过的每一层都是简单模 \(\kappa\)，共有 \(n\) 层。写 \(f=u\pi^n\) 后 \(fR=\pi^nR\)；单位情形 \(n=0\) 的商模为零，长度同样为零。同样的乘法同构还给出 \(\mathfrak m/\mathfrak m^2\cong\kappa\)。最后，\(t\in\mathfrak m\setminus\mathfrak m^2\) 等价于其阶为 \(1\)，也就等价于它生成 \(\mathfrak m\)。
\end{proof}

\subsection{Dedekind 整环的局部观点}
\label{g5:subsec:0401:03}

一个环可以在每个非零素理想处都呈现上述单参数结构，却没有一个元素同时生成所有理想。局部的阶因此应当附着在各个素理想上；把有限多个非零阶放在一起，就得到整体理想的分解。

\begin{definition}
\label{g5:def:dedekind-domain}
	一个 Noether 整闭整环 \(A\)，若 \(\dim A\leq1\)，就称为\term[Dedekindzhenghuan]{Dedekind 整环}[Dedekind domain]。本节采用允许域的约定；非域的 Dedekind 整环恰为一维 Noether 整闭整环。
\end{definition}

Dedekind 整环用理想的乘积分解取代元素的唯一分解。\footnote{Julius Wilhelm Richard Dedekind（1831--1916），德国数学家，发展了代数整数的理想理论，并用实数分割刻画无理数。}先说明怎样从各点的局部信息恢复环与理想，这也将解释整闭性为何可以逐点检验。

\begin{proposition}
\label{g5:prop:domain-local-intersection}
	设 \(A\) 为整环，\(K=\operatorname{Frac}(A)\)。对任意 \(A\) 子模 \(M\subseteq K\)，在 \(K\) 中有
	\[
		M=\bigcap_{\mathfrak m\in\operatorname{Max}A}M_{\mathfrak m}.
	\]
	特别地，\(A=\bigcap_{\mathfrak m}A_{\mathfrak m}\)，两个分式理想相等当且仅当它们在每个极大理想处的局部化相等。而且 \(A\) 整闭，当且仅当每个 \(A_{\mathfrak m}\) 都整闭。
\end{proposition}
\begin{proof}
	各个局部化可视为 \(K\) 的子模，故从左到右的包含成立。若 \(z\in K\setminus M\)，理想
	\[
		J=\{a\in A:az\in M\}
	\]
	不含 \(1\)，可放入某个极大理想 \(\mathfrak m\)。若 \(z\in M_{\mathfrak m}\)，就存在 \(s\notin\mathfrak m\) 使 \(sz\in M\)，与 \(s\in J\subseteq\mathfrak m\) 矛盾。所以 \(z\) 不属于右侧的交，证明等式。

	若 \(A\) 整闭，则\cref{g5:prop:integral-localization-quotient}保证它的局部化整闭。反过来，若每个 \(A_{\mathfrak m}\) 整闭，而 \(z\in K\) 整于 \(A\)，同一首一方程表明 \(z\) 整于每个 \(A_{\mathfrak m}\)，故 \(z\in A_{\mathfrak m}\) 对所有 \(\mathfrak m\) 成立。交等式于是给出 \(z\in A\)。
\end{proof}

\begin{theorem}
\label{g5:thm:dedekind-local}
	设 \(A\) 为 Noether 整环。则 \(A\) 为 Dedekind 整环，当且仅当对每个非零素理想 \(\mathfrak p\)，局部环 \(A_{\mathfrak p}\) 都为 DVR。
\end{theorem}
\begin{proof}
	若 \(A\) 为 Dedekind 整环，则每个非零素理想的高度为 \(1\)。局部化 \(A_{\mathfrak p}\) 因而是一维 Noether 局部整环，且仍整闭，由\cref{g5:thm:dvr-characterizations}知它为 DVR。

	反过来，若 \(A\) 为域，结论已经成立。否则每个极大理想均非零，按假设它的局部环为 DVR，所以每个极大理想的高度都为 \(1\)。任何有限素链都可放在某个极大理想下方，故 \(\dim A=1\)。每个极大理想处的局部环又整闭，由\cref{g5:prop:domain-local-intersection}得 \(A\) 整闭，因此 \(A\) 为 Dedekind 整环。
\end{proof}

对非零素理想 \(\mathfrak p\)，记 \(A_{\mathfrak p}\) 的规范化阶为 \(v_{\mathfrak p}\)。若 \(I\) 是非零理想，局部化后有
\[
	I A_{\mathfrak p}
	=(\mathfrak pA_{\mathfrak p})^{n_{\mathfrak p}},
	\qquad n_{\mathfrak p}\geq0.
\]
整数 \(n_{\mathfrak p}\) 描述 \(I\) 在这个点的全部信息。它可以不同于 \(I\) 的整体生成元数，也不要求存在一个统一的生成元在所有点同时实现这些阶。

\begin{theorem}
\label{g5:thm:dedekind-factorization}
	设 \(A\) 为 Dedekind 整环，\(K=\operatorname{Frac}(A)\)。
	\begin{enumerate}
		\item 每个非零理想 \(I\subseteq A\) 都唯一写为
		\[
			I=\mathfrak p_1^{n_1}\cdots\mathfrak p_r^{n_r},
		\]
		其中 \(\mathfrak p_i\) 为互不相同的非零素理想，\(n_i\geq1\)。\(I=A\) 对应空乘积。
		\item 每个分式理想均可逆，并唯一写为有限乘积
		\[
			I=\prod_{\mathfrak p\ne(0)}\mathfrak p^{n_{\mathfrak p}},
			\qquad n_{\mathfrak p}\in\mathbb Z,
		\]
		其中只有有限多个指数非零，负次幂由可逆理想的逆定义。每个指数由 \(IA_{\mathfrak p}\) 的阶确定。
	\end{enumerate}
\end{theorem}
\begin{proof}
	域的非零理想只有自身，结论成立。以下设 \(A\) 不是域。先取非零真理想 \(I\)。包含 \(I\) 的极大理想只有有限多个：选 \(0\ne a\in I\)，则 \(A/(a)\) 是 Noether 环；其素理想对应于 \(A\) 中包含非零元素 \(a\) 的素理想，故全部极大。因此这个商环零维，由\cref{g5:thm:artinian-zero-dimension}，它只有有限多个素理想。包含 \(I\) 的极大理想也就只有有限多个，记为 \(\mathfrak p_1,\ldots,\mathfrak p_r\)。

	对每个 \(\mathfrak p_i\)，由 DVR 中的理想刻画，有
	\[
		IA_{\mathfrak p_i}
		=(\mathfrak p_iA_{\mathfrak p_i})^{n_i},
		\qquad n_i\geq1.
	\]
	令 \(J=\prod_i\mathfrak p_i^{n_i}\)。不同极大理想 \(\mathfrak p_i,\mathfrak q\) 之间可选元素 \(s\in\mathfrak p_i\setminus\mathfrak q\)，所以 \(\mathfrak p_iA_{\mathfrak q}=A_{\mathfrak q}\)。因此在每个极大理想处，\(J\) 的局部化都与 \(I\) 相同：在 \(\mathfrak p_i\) 处只保留对应的一项，在其他点处两者都成为单位理想。由\cref{g5:prop:domain-local-intersection}，\(I=J\)。再次局部化到各个 \(\mathfrak p_i\)，DVR 中指数的唯一性给出整体分解的唯一性。

	下面证明分式理想 \(I\) 可逆。因为 \(A\) Noether，\(I\) 有有限生成元 \(a_1,\ldots,a_s\)。对每个极大理想 \(\mathfrak m\)，有
	\[
		(I^{-1})_{\mathfrak m}=(I_{\mathfrak m})^{-1},
	\]
	右侧的逆在 \(A_{\mathfrak m}\) 中计算。一个包含由定义得到。反过来，若 \(xI_{\mathfrak m}\subseteq A_{\mathfrak m}\)，则对每个 \(a_i\) 可选 \(s_i\notin\mathfrak m\)，使 \(s_i x a_i\in A\)。取 \(s=\prod_i s_i\)，就有 \(sxI\subseteq A\)，即 \(sx\in I^{-1}\)，从而 \(x\in(I^{-1})_{\mathfrak m}\)。这里有限生成保证能用同一个分母处理所有生成元。

	在 DVR \(A_{\mathfrak m}\) 中，分式理想为主理想，所以
	\[
		(II^{-1})_{\mathfrak m}
		=I_{\mathfrak m}(I_{\mathfrak m})^{-1}
		=A_{\mathfrak m}.
	\]
	逐点的等式由交公式提升为 \(II^{-1}=A\)，证明可逆性。

	最后选 \(0\ne d\in A\) 使 \(dI\subseteq A\)。将普通理想 \(dI\) 与 \(dA\) 分别作第一项的分解，并在可逆分式理想的乘法下相除，就得到 \(I\) 的整数指数分解。其局部指数是 \(dI\) 与 \(dA\) 的指数差；再次局部化即可证明唯一性。
\end{proof}

这把各点的单参数结构拼成了整体理想分解；与之对应的一般理想可逆性判据和 Dedekind 刻画见\cite[第11节，定理11.3、11.6，第80--83页]{Matsumura1989}。每个分式理想还是一个秩一有限投射模：若 \(1=\sum_i a_ib_i\)、\(a_i\in I,b_i\in I^{-1}\)，则映射 \(\lambda_i(x)=b_ix\) 取值于 \(A\)，并满足
\[
	x=\sum_i\lambda_i(x)a_i.
\]
由\cref{g5:prop:projective-dual-basis}，这就是有限投射性的对偶基表达。局部化到分式域后 \(I\otimes_AK=K\)，所以秩为一。整体主理想相当于这个投射模有一组含一个元素的自由基；可逆性本身并不保证它存在。

\begin{example}
\label{g5:exm:dedekind-nonprincipal}
	令 \(A=\mathbb Z[\sqrt{-5}]\subset\mathbb C\)。它是 Dedekind 整环，却不是主理想整环。先核对整闭性。若 \(z=r+s\sqrt{-5}\in\mathbb Q(\sqrt{-5})\)、\(r,s\in\mathbb Q\)，且 \(z\) 整于 \(A\)，则因 \(A\) 有限于 \(\mathbb Z\)，整性传递说明 \(z\) 整于 \(\mathbb Z\)。若 \(s=0\)，\(\mathbb Z\) 的整闭性给出 \(z\in\mathbb Z\)。若 \(s\ne0\)，最小多项式为
	\[
		T^2-2rT+(r^2+5s^2).
	\]
	由\cref{g5:lem:normal-minimal-polynomial}，\(2r\) 与 \(r^2+5s^2\) 都为整数。因此 \(20s^2\) 为整数。将 \(s\) 写成既约分数，其分母的平方整除 \(20\)，所以 \(2s\) 也是整数。令 \(a=2r,b=2s\)，则
	\[
		a^2+5b^2\equiv0\pmod4.
	\]
	平方模 \(4\) 只有 \(0,1\)，这迫使 \(a,b\) 均为偶数。因此 \(r,s\in\mathbb Z\)，即 \(z\in A\)，证明整闭。又 \(A\) 是有限 \(\mathbb Z\) 模，故 Noether；整扩张保持维数，所以 \(\dim A=1\)。

	理想 \(\mathfrak p=(2,1+\sqrt{-5})\) 的商环为 \(\mathbb F_2\)，因而是极大理想。若 \(\mathfrak p=\alpha A\)，则 \(\alpha\) 同时整除 \(2\) 和 \(1+\sqrt{-5}\)。取范数
	\[
		N(a+b\sqrt{-5})=a^2+5b^2
	\]
	便知正整数 \(N(\alpha)\) 同时整除 \(4\) 和 \(6\)。因为 \(\mathfrak p\) 为真理想，\(\alpha\) 不是单位，而范数为 \(1\) 的元素可用共轭求逆，故只能有 \(N(\alpha)=2\)。但整数方程 \(a^2+5b^2=2\) 无解，矛盾。另一方面，直接计算得
	\[
		\mathfrak p^2=(2).
	\]
	确实，平方理想的三个生成元 \(4,2+2\sqrt{-5},-4+2\sqrt{-5}\) 都被 \(2\) 整除；后两个的差为 \(6\)，与 \(4\) 一起生成 \(2\)。因此 \(\mathfrak p^{-1}=\tfrac12\mathfrak p\)。这个非主理想在每个局部环中都有一个生成元，整体上仍需多个生成元。
\end{example}

\subsection{正规曲线的局部参数}
\label{g5:subsec:0401:04}

现在把阶函数放回曲线。设 \(k\) 为代数闭域，\(X\) 为不可约仿射代数集，且 \(\dim k[X]=1\)。若坐标环 \(k[X]\) 整闭，就称 \(X\) 为\term[zhengguiquxian]{正规曲线}[normal curve]。由局部交公式，这等价于每个点 \(p\in X\) 的局部环 \(\mathcal O_{X,p}\) 整闭。曲线的函数域记为 \(k(X)=\operatorname{Frac}k[X]\)，它也是每个局部环的分式域。

\begin{proposition}
\label{g5:prop:normal-curve-parameters}
	设 \(k\) 为代数闭域，\(X\) 为正规不可约仿射曲线，\(p\in X\)。则 \(\mathcal O_{X,p}\) 为 DVR，剩余域为 \(k\)。存在正则函数 \(t\in k[X]\)，使 \(t(p)=0\) 且 \(t\) 的芽生成局部环的极大理想；这样的芽称为 \(p\) 处的\term[jubucanshu]{局部参数}[local parameter]。每个 \(f\in k(X)^\times\) 唯一写为
	\[
		f=u\,t^n,\qquad u\in\mathcal O_{X,p}^{\times},\quad n\in\mathbb Z.
	\]
	指数 \(n\) 与局部参数的选择无关，记为 \(\operatorname{ord}_p(f)\)。此外：
	\begin{enumerate}
		\item \(f\) 在 \(p\) 处正则，当且仅当 \(\operatorname{ord}_p(f)\geq0\)；
		\item \(f\) 在 \(p\) 处正则且取值为零，当且仅当 \(\operatorname{ord}_p(f)>0\)；
		\item 若 \(0\ne f\in\mathcal O_{X,p}\)，则
		\[
			\ell_{\mathcal O_{X,p}}\bigl(\mathcal O_{X,p}/f\mathcal O_{X,p}\bigr)
			=\operatorname{ord}_p(f).
		\]
	\end{enumerate}
\end{proposition}
\begin{proof}
	令 \(A=k[X]\)，它是 Noether 一维整闭整环。点 \(p\) 对应非零极大理想 \(\mathfrak m_p\)，且 \(\mathcal O_{X,p}\cong A_{\mathfrak m_p}\)，因此\cref{g5:thm:dedekind-local}给出 DVR 性。由弱零点定理，剩余域为 \(k\)。

	在局部环中取统一化元 \(\pi=a/s\)，其中 \(a,s\in A\)、\(s(p)\ne0\)。因为 \(s\) 在局部环中为单位，\(a=s\pi\) 也是统一化元；于是可取全局正则函数 \(t=a\)，且 \(t(p)=0\)。其余结论由 DVR 的唯一指数表达、单位与极大理想的阶判据，以及\cref{g5:prop:dvr-length}得到。第一项还使用了第一章的函数芽描述：有理函数在 \(p\) 处正则，恰指它属于 \(\mathcal O_{X,p}\)。
\end{proof}

若 \(\operatorname{ord}_p(f)=n>0\)，就说 \(f\) 在 \(p\) 处有 \(n\) 阶零点；若 \(\operatorname{ord}_p(f)=-n<0\)，就说它有 \(n\) 阶极点。对正规仿射曲线，非零有理函数只有有限多个零点和极点。事实上，Dedekind 分解将主分式理想写成
\[
	f\,k[X]=\prod_{p\in X}\mathfrak m_p^{\operatorname{ord}_p(f)},
\]
其中只有有限多个指数非零；局部化到 \(\mathfrak m_p\) 后，指数恰是局部环中的阶。在仿射直线上，\(f=t^2/(t-1)^3\) 的阶在 \(0\) 处为 \(2\)，在 \(1\) 处为 \(-3\)，在其他点为零。

\ProseParagraphBreak

局部参数只要求在指定点消失一次，不要求它生成整条曲线的函数域。一个方程算例可以看清这个区别。

\begin{example}
\label{g5:exm:local-parameter-not-field-generator}
	设 \(k\) 为特征不等于 \(2\) 的代数闭域，考虑曲线
	\[
		X:\ y^2=x^3+1,\qquad p=(0,1).
	\]
	多项式 \(x^3+1\) 在 \(k(x)\) 中不是平方：一个既约有理函数的平方若是多项式，分母必须为常数，而平方多项式的次数为偶数。因此 \(y^2-x^3-1\) 不可约，坐标环为整环。

	在点 \(p\) 的局部环中，\(y+1\) 的值为 \(2\)，故为单位。由
	\[
		(y-1)(y+1)=x^3
	\]
	得到 \(y-1=x^3/(y+1)\)，所以极大理想满足
	\[
		(x,y-1)\mathcal O_{X,p}=x\mathcal O_{X,p}.
	\]
	这是非域 Noether 局部整环的非零主极大理想，故由\cref{g5:thm:dvr-characterizations}，\(\mathcal O_{X,p}\) 为 DVR，且 \(x\) 为局部参数。于是
	\[
		\operatorname{ord}_p(x)=1,\qquad
		\operatorname{ord}_p(y-1)=3.
	\]
	同时 \(k(X)/k(x)\) 为二次扩张，所以这个局部参数并不单独生成函数域。这里局部环的阶直接从方程与一个可逆因子求出。
\end{example}

一维整闭性的局部结论最终落在统一化元上：它给出点附近唯一的阶、全部非零分式理想及主理想商的长度。把这些局部环放回一个坐标环时，Dedekind 分解负责记录各点的指数。下一节将构造整闭包，使原曲线缺失的整函数也进入这一局部结构。

% 纲要条目排为编号小节标题，并进入目录与书签；说明段仍保留于纲要注释中。

% !TeX root = ../../Book5.tex
% 纲要标题：### 4.2 正规化的构造与有限性
% 纲要位置：几何学计划纲要.md，第 5158 行。
% 正文状态：初稿；尚未正式审核。

% 纲要内容开始
% 
% * 分式域及总分式环中的整闭包
% * Noether 正规化后的有限性归约
% * 有限可分扩张的迹配对与分母控制
% * 不可分扩张与整闭包有限性
% 
%   只涉及有限个系数的纯不可分扩张，不能默设基域完美或 Frobenius 有限。
% 
% * 约化对象的各分支与粘合
% 
% 纲要内容结束

\section{正规化的构造与有限性}
\label{g5:sec:0402}

第三章已经证明，尖点的坐标环 \(k[t^2,t^3]\) 在分式域中缺少一个整元 \(t\)。补入这个元素便得到 \(k[t]\)，每个闭点附近也随之得到离散赋值环。一般环可能需要补入许多整元；即使每个元素分别满足首一方程，也不能立即断言全部整元能由有限个元素共同生成。本节先确定补入元素的环境，再证明有限型代数的整闭包确实是有限模。

\ProseParagraphBreak

不可约对象只有一个函数域，可约对象则需要同时保留各分支的函数域。处理有限性时，Noether 正规化把基环化为多项式环；可分扩张中的迹配对提供统一分母，不可分扩张则通过有限个系数的开根进入可分情形。这两步分别解决不同的障碍。

\subsection{分式域及总分式环中的整闭包}
\label{g5:subsec:0402:01}

\begin{definition}\label{g5:def:normalization-ring}
	设 \(A\) 为整环，或为约化 Noether 环，\(S\) 为 \(A\) 中全部非零因子组成的乘法集。环
	\[
		Q(A)=S^{-1}A
	\]
	称为 \(A\) 的\term[zongfenshihuan]{总分式环}[total ring of fractions]。\(A\) 在 \(Q(A)\) 中的整闭包
	\[
		A^\nu=\{z\in Q(A):z\text{ 在 }A\text{ 上整}\}
	\]
	称为 \(A\) 的\term[zhengguihua]{正规化}[normalization]。若 \(A\) 为整环，则 \(Q(A)=\operatorname{Frac}A\)，这就是上一章所说的普通正规化。
\end{definition}
\symidx{\(Q(A)\)：约化环的总分式环}
\symidx{\(A^\nu\)：环的正规化}

非零因子作为分母使 \(A\to Q(A)\) 保持单射。若任意反演零因子，就可能删去一个分支。例如在 \(k[x,y]/(xy)\) 中反演 \(x\) 会使 \(y=0\)，于是只保留一条轴。总分式环的目标是保留全部分支，再分别形成函数域。

\begin{proposition}\label{g5:prop:total-quotient-components}
	设 \(A\ne0\) 为约化 Noether 环，\(\mathfrak p_1,\ldots,\mathfrak p_r\) 为全部极小素理想，记
	\[
		A_i=A/\mathfrak p_i,\qquad K_i=\operatorname{Frac}(A_i).
	\]
	则自然同态给出
	\[
		Q(A)\cong\prod_{i=1}^r K_i,
		\qquad
		A^\nu\cong\prod_{i=1}^r A_i^\nu.
	\]
\end{proposition}
\begin{proof}
	极小素理想有限且交为幂零根；\(A\) 约化，所以 \(\bigcap_i\mathfrak p_i=0\)。因此 \(A\to\prod_i A_i\) 单射。对每个 \(i\)，从各 \(\mathfrak p_j\setminus\mathfrak p_i\)、\(j\ne i\)，取一个元素并相乘，得到
	\[
		b_i\in\bigcap_{j\ne i}\mathfrak p_j,\qquad b_i\notin\mathfrak p_i.
	\]
	当 \(r=1\) 时取 \(b_1=1\)。这些元素在乘积中只在第 \(i\) 个分量非零。

	一个元素 \(a\in A\) 为非零因子，当且仅当它不属于任何 \(\mathfrak p_i\)。若各分量均非零，\(ac=0\) 在每个整环 \(A_i\) 中迫使 \(c=0\)，得到一向。若 \(a\in\mathfrak p_i\)，则 \(ab_i=0\)，而 \(b_i\ne0\)，得到另一向。因此局部化后有单射 \(Q(A)\to\prod_i K_i\)。

	下面证明满射，而不只把两侧看成相似的分式表达。令 \(s=\sum_i b_i\)。它模每个 \(\mathfrak p_i\) 等于非零的 \(b_i\)，故为非零因子。元素 \(b_i/s\) 就是第 \(i\) 个坐标幂等元。更一般地，给定 \(a,c\in A\)、\(c\notin\mathfrak p_i\)，令
	\[
		d=cb_i+\sum_{j\ne i}b_j.
	\]
	它在每个分量都非零，因而也是合法分母。分式 \(ab_i/d\) 的第 \(i\) 个分量为 \(\overline a/\overline c\)，其余分量为零。将这些分式相加，便得到任意 \((z_1,\ldots,z_r)\in\prod_i K_i\)，证明第一式。

	若 \(z\in Q(A)\) 整于 \(A\)，它的每个分量显然整于 \(A_i\)。反过来，设每个 \(z_i\) 整于 \(A_i\)，为它取一个首一方程，并把系数提升为 \(A\) 的元素，得到 \(f_i(T)\in A[T]\)。乘积
	\[
		f(T)=\prod_{i=1}^r f_i(T)
	\]
	仍首一；在第 \(i\) 个分量中有 \(f_i(z_i)=0\)，所以 \(f(z)=0\) 于整个乘积环。这证明 \(z\) 整于 \(A\)，得到第二式。
\end{proof}

这个分解说明，可约对象的正规化会把不同不可约分支分开。乘积环中的坐标幂等元满足 \(e_i^2-e_i=0\)，原环中若没有这些幂等元，正规化就会把它们补入。分支内部仍可能有进一步的整元需要补入，如尖点的参数 \(t\)。

\ProseParagraphBreak

对一般环，称每个素理想处的局部环都是整闭整环的环为\term[zhengguihuan]{正规环}[normal ring]。这里局部环为整环也是条件的一部分。对于约化 Noether 环，这个定义与 \(A=A^\nu\) 相容；一个非平凡乘积环本身不是整环，却仍然可以正规。

\begin{proposition}\label{g5:prop:normalization-localization}
	设 \(A\) 为约化 Noether 环，\(S\subseteq A\) 为乘法集。则自然识别给出
	\[
		S^{-1}(A^\nu)\cong(S^{-1}A)^\nu.
	\]
	此外，\(A^\nu\) 为正规环，并且 \(A\) 正规当且仅当 \(A=A^\nu\)。
\end{proposition}
\begin{proof}
	在\cref{g5:prop:total-quotient-components}的乘积表示中，把 \(S\) 反演会删去那些 \(\mathfrak p_i\) 与 \(S\) 相交的分量，其余分量仍是同一个域 \(K_i\)。极小素理想的局部化对应因此给出
	\[
		S^{-1}Q(A)\cong Q(S^{-1}A).
	\]
	第三章的整闭包局部化公式\cref{g5:prop:integral-localization-quotient}随即给出第一式；若局部化为零环，两边也同为零环。

	每个 \(A_i^\nu\) 在 \(K_i\) 中整闭：一个在 \(A_i^\nu\) 上整的元素，由整性传递仍在 \(A_i\) 上整，因而已属于 \(A_i^\nu\)。整闭整环的局部化仍整闭。有限乘积环的每个素局部环来自其中一个因子，所以乘积 \(A^\nu=\prod_i A_i^\nu\) 为正规环。特别地，\(A=A^\nu\) 蕴含 \(A\) 正规。

	反过来，设所有 \(A_{\mathfrak m}\)、\(\mathfrak m\) 为极大理想，都为整闭整环。局部化公式说明
	\[
		(A^\nu)_{\mathfrak m}=(A_{\mathfrak m})^\nu=A_{\mathfrak m}.
	\]
	故模 \(A^\nu/A\) 在每个极大理想处局部化为零。由第二章的模消失局部检测原则，它本身为零，即 \(A=A^\nu\)。这个检测原则对任意模成立，尚未使用正规化的有限性。
\end{proof}

构造还具有一个实用的唯一性。若 \(A\to C\) 是整环之间的单射，且 \(C\) 整闭，则它唯一延伸为 \(A^\nu\to C\)：把 \(\operatorname{Frac}A\) 嵌入 \(\operatorname{Frac}C\)，每个 \(z\in A^\nu\) 都整于 \(C\)，因而落在 \(C\) 中；分式表示则强制了像的唯一性。这说明，从正规对象到原对象的支配映射可以经过正规化来表达。

\subsection{Noether 正规化后的有限性归约}
\label{g5:subsec:0402:02}

要证明正规化有限，不能只逐个收集整性方程。我们需要找到一个有限模，把全部待加入的元素同时装进去；Noether 性随后保证其中的子模仍有限生成。先把这个任务归约到多项式基环。

\begin{proposition}\label{g5:prop:normalization-polynomial-reduction}
	设 \(k\) 为域，\(A\) 为有限型 \(k\) 整代数，\(F=\operatorname{Frac}A\)，\(L/F\) 为有限域扩张。取 Noether 正规化
	\[
		C=k[x_1,\ldots,x_d]\subseteq A,
	\]
	其中 \(x_i\) 代数独立，\(A\) 是有限 \(C\) 模。则 \(L\) 是 \(K=\operatorname{Frac}C\) 的有限扩张，且 \(A\) 在 \(L\) 中的整闭包等于 \(C\) 在 \(L\) 中的整闭包。若后者为有限 \(C\) 模，它也为有限 \(A\) 模。
\end{proposition}
\begin{proof}
	第三章正规化的超越次数证明已经表明：把 \(C\setminus\{0\}\) 反演后，\(A\) 成为有限维 \(K\) 整代数，因而成为域；它恰是 \(F\)。所以 \(F/K\) 有限，再用域扩张次数的乘法得 \(L/K\) 有限。

	因 \(A\) 整于 \(C\)，一个整于 \(A\) 的元素由传递性也整于 \(C\)；反方向只需把同一个首一方程的系数从 \(C\) 看作 \(A\) 的元素。因此两份整闭包相同，记为 \(D\)。若 \(d_1,\ldots,d_s\) 生成 \(D\) 作为 \(C\) 模，则
	\[
		D=\sum_i Cd_i\subseteq\sum_i Ad_i\subseteq D,
	\]
	最后一个包含来自 \(D\) 为含 \(A\) 的环。所以同一组元素也生成 \(D\) 作为 \(A\) 模。
\end{proof}

多项式环 \(C\) 同时具有 Noether 性与整闭性。前者来自 Hilbert 基定理。后者可由唯一分解说明，下面回顾所需的代数依据。

\begin{lemma}\label{g5:lem:polynomial-ring-normal}
	域 \(k\) 上的有限变量多项式环 \(k[x_1,\ldots,x_d]\) 是唯一分解整环，因而整闭。
\end{lemma}
\begin{proof}
	只需说明唯一分解整环 \(R\) 的多项式环 \(R[T]\) 仍唯一分解，再从域出发归纳。称系数没有共同不可约因子的多项式为本原多项式。两个本原多项式的乘积仍本原：若某个不可约元 \(p\in R\) 整除乘积的全部系数，模 \(p\) 后会得到两个非零多项式的乘积为零，而 \(R/(p)\) 是整环，矛盾。

	每个非零多项式可唯一写成系数公因子乘本原多项式，公因子仅差单位。令 \(E=\operatorname{Frac}R\)。把 \(E[T]\) 中分解的各因子清去分母，再除去系数公因子，便得到 \(R[T]\) 中的本原因子。刚证的乘积性质保证所得等式中的剩余标量为 \(R\) 的单位：若两个本原多项式相差 \(a/b\in E^\times\)，把 \(a,b\) 约为互素后，系数比较迫使 \(a,b\) 都为单位。因此 \(E[T]\) 的因子分解给出 \(R[T]\) 的分解；唯一性分别来自 \(R\) 的系数公因子分解与 \(E[T]\) 的多项式分解。于是 \(R[T]\) 唯一分解。最后应用\cref{g5:prop:ufd-integrally-closed}得到整闭性。
\end{proof}

接下来先在一般整闭 Noether 整环上处理有限可分扩张，再利用多项式的有限系数特征处理不可分扩张。这个归约及两种情形的组织可对照\cite[第33节，整闭包讨论及引理1，第262--263页]{Matsumura1989}。

\subsection{有限可分扩张的迹配对与分母控制}
\label{g5:subsec:0402:03}

有限域扩张 \(L/K\) 中，每个 \(z\in L\) 给出 \(K\) 线性映射 \(m_z:L\to L\)、\(u\mapsto zu\)。定义
\[
	\operatorname{Tr}_{L/K}(z)=\operatorname{tr}(m_z).
\]
这是域扩张的\term[yukuozhangdeji]{迹}[field trace]；相应双线性形式
\[
	(z,w)\longmapsto\operatorname{Tr}_{L/K}(zw)
\]
称为\term[jipeidui]{迹配对}[trace pairing]。矩阵迹不依赖所选基，所以这两个定义也不依赖坐标。
\symidx{\(\operatorname{Tr}_{L/K}\)：有限域扩张的迹}

\ProseParagraphBreak

一个代数元可分，是指它的最小多项式在分裂域中没有重根；有限扩张可分，是指其中每个元素都可分。可分性使迹配对非退化，从而允许把一个元素的坐标恢复为若干个迹的线性组合。下面把这个线性代数事实及其所需的域论说明出来。

\begin{lemma}\label{g5:lem:separable-trace-pairing}
	设 \(L/K\) 为次数 \(n\) 的有限可分域扩张。则迹配对非退化。若 \(R\subseteq K\) 为整闭整环且 \(K=\operatorname{Frac}R\)，\(z\in L\) 整于 \(R\)，则
	\[
		\operatorname{Tr}_{L/K}(z)\in R.
	\]
\end{lemma}
\begin{proof}
	取包含 \(L\) 的代数闭域 \(\Omega\)。可分性给出
	\[
		L\otimes_K\Omega\cong\prod_{\sigma:L\hookrightarrow\Omega}\Omega
		\cong\Omega^n,\qquad z\otimes c\longmapsto(\sigma(z)c)_\sigma.
	\]
	这个同构可以按有限组代数生成元逐步验证。添入一个生成元 \(\alpha\) 时，其最小多项式可分，沿已有的每个嵌入在 \(\Omega\) 上分解为互异一次因子；中国剩余定理把相应商代数分成这些根对应的 \(\Omega\) 因子。重复操作后，因子数等于各步扩张次数的乘积 \(n\)，也恰好列出全部嵌入。这里没有要求 \(L/K\) 为正规扩张。

	在这个乘积中，乘以 \(z\) 的线性映射为对角映射，对角元为 \(\sigma(z)\)。迹与标量扩张相容，因为它直接由矩阵对角元相加得到。因此
	\[
		\operatorname{Tr}_{L/K}(z)=\sum_\sigma\sigma(z).
	\]
	扩张标量后的迹配对成为 \(\Omega^n\) 上的
	\((a,b)\mapsto\sum_\sigma a_\sigma b_\sigma\)，显然非退化。若原迹矩阵的行列式为零，扩张标量后仍为零，矛盾。因此原配对也非退化。

	若 \(z\) 整于 \(R\)，每个 \(\sigma(z)\) 都满足同一个 \(R\) 系数首一方程，故都整于 \(R\)。整元组成子环，所以它们的和整于 \(R\)。这个和又是属于 \(K\) 的迹值，\(R\) 整闭便迫使它属于 \(R\)。
\end{proof}

\begin{theorem}\label{g5:thm:separable-integral-closure-finite}
	设 \(R\) 为整闭 Noether 整环，\(K=\operatorname{Frac}R\)，\(L/K\) 为有限可分扩张。则 \(R\) 在 \(L\) 中的整闭包 \(D\) 是有限 \(R\) 模。
\end{theorem}
\begin{proof}
	先选一组由整元组成的 \(K\) 基。对任意代数元 \(b\in L\)，将其首一 \(K\) 系数方程的分母合并为 \(0\ne a\in R\)，便知 \(ab\) 满足一个首一 \(R\) 系数方程：若
	\[
		b^m+c_1b^{m-1}+\cdots+c_m=0,
	\]
	且 \(ac_j\in R\)，则 \(ab\) 的方程系数为 \(a^jc_j\in R\)。因此可以将任一 \(K\) 基分别乘以非零 \(R\) 元素，得到
	\[
		e_1,\ldots,e_n\in D.
	\]

	令 \(M=(\operatorname{Tr}_{L/K}(e_ie_j))_{i,j}\)，并记 \(\Delta=\det M\)。由上个引理，\(M\) 的各项属于 \(R\)，且 \(\Delta\ne0\)。任取 \(z\in D\)，写
	\[
		z=\sum_i c_i e_i,\qquad c_i\in K.
	\]
	对每个 \(j\) 取迹，得到
	\[
		\sum_i c_i\operatorname{Tr}_{L/K}(e_ie_j)
		=\operatorname{Tr}_{L/K}(ze_j)\in R.
	\]
	右侧属于 \(R\)，因为 \(ze_j\) 仍整。用伴随矩阵解这个方程组，得到 \(\Delta c_i\in R\)。因此
	\begin{equation}\label{g5:eq:trace-common-denominator}
		D\subseteq\Delta^{-1}(Re_1+\cdots+Re_n).
	\end{equation}
	右侧是有限自由 \(R\) 模。左侧是其中的 \(R\) 子模，而 \(R\) Noether，所以 \(D\) 有限生成。
\end{proof}

同一个非零元素 \(\Delta\) 控制了所有整元的分母，这正是逐个整性方程未能直接提供的信息。该定理是\cite[第33节，引理1，第262页]{Matsumura1989}的有限可分版本；迹配对写法把分母控制明确化为一个可逆矩阵。

\ProseParagraphBreak

不可分扩张中迹配对可能完全失效。例如在特征 \(p>0\) 下，若 \(L=K(\alpha)\)、\(\alpha^p=a\) 且 \([L:K]=p\)，则基 \(1,\alpha,\ldots,\alpha^{p-1}\) 下，乘以 \(\alpha^j\)、\(1\leq j<p\)，没有对角项；乘以 \(1\) 的迹又是 \(p=0\)。线性性因此给出 \(\operatorname{Tr}_{L/K}\equiv0\)。这时应改变基域以分离不可分部分，而不能继续使用非零的迹行列式。

\subsection{不可分扩张与整闭包有限性}
\label{g5:subsec:0402:04}

以下设特征为 \(p>0\)。若有限扩张 \(E/K\) 的每个元素的某个 \(p\) 次幂迭代落回 \(K\)，就称它为\term[chunbukefenkuozhang]{纯不可分扩张}[purely inseparable extension]。有限个域生成元可取共同指数 \(q=p^e\)，使每个生成元的 \(q\) 次幂都属于 \(K\)。在代数闭域中 \(q\) 次根唯一，因为 \(u^q-v^q=(u-v)^q\)。

\begin{lemma}\label{g5:lem:finite-coefficient-root-extension}
	设 \(k\) 为特征 \(p>0\) 的域，
	\[
		C=k[x_1,\ldots,x_d],\qquad K=k(x_1,\ldots,x_d),
	\]
	其中 \(x_i\) 代数独立，\(L/K\) 为有限域扩张。则存在 \(q=p^e\)、有限纯不可分扩张
	\[
		k'=k(a_1^{1/q},\ldots,a_s^{1/q}),\qquad a_j\in k,
	\]
	满足如下性质：记
	\[
		z_i=x_i^{1/q},\quad C'=k'[z_1,\ldots,z_d],
		\quad K'=\operatorname{Frac}C',
	\]
	则 \(C'\) 是整闭 Noether 整环且为有限 \(C\) 模，合成域 \(LK'\) 为 \(K'\) 的有限可分扩张。
\end{lemma}
\begin{proof}
	在一个固定代数闭域中进行全部开根和合成。取有限组域生成元 \(L=K(\alpha_1,\ldots,\alpha_m)\)。设 \(\alpha_j\) 的首一最小多项式为 \(f_j(T)\)。在特征 \(p\) 下，导数为零恰表示所有出现的指数都被 \(p\) 整除。反复除去这个共同指数，可写
	\[
		f_j(T)=g_j(T^{q_j}),\qquad q_j=p^{e_j},
	\]
	其中 \(g_j\in K[T]\) 首一、不可约且导数非零。不可约性来自：若 \(g_j\) 可分解，代入 \(T^{q_j}\) 就分解了 \(f_j\)。因此 \(g_j\) 无重根。

	所有 \(g_j\) 的系数都是 \(x_i\) 的有理函数。选择分子、分母的多项式表达，并把其中出现的全部 \(k\) 系数列为 \(a_1,\ldots,a_s\)。这是有限列表。取 \(q=p^e\) 同时为所有 \(q_j\) 的倍数，并构造题述 \(k'\)、\(C'\)、\(K'\)。对任一上述有理函数 \(c=P(x)/Q(x)\)，把系数和变量分别开 \(q\) 次根便得到
	\[
		c^{1/q}
		=\frac{P^{[1/q]}(z_1,\ldots,z_d)}
		{Q^{[1/q]}(z_1,\ldots,z_d)}\in K',
	\]
	其中上标 \([1/q]\) 表示对系数逐项开根。分母非零，因为它的 \(q\) 次幂为原非零多项式。由 \(q_j\mid q\)，各系数的 \(q_j\) 次根也属于 \(K'\)。

	将 \(g_j\) 的系数逐项开 \(q_j\) 次根，得到 \(h_j(T)\in K'[T]\)。于是
	\[
		h_j(T)^{q_j}=g_j(T^{q_j})=f_j(T),
		\qquad h_j(\alpha_j)=0.
	\]
	\(h_j\) 无重根。确实，在代数闭域中若
	\(g_j(T)=\prod_\ell(T-\beta_\ell)\)，其根 \(\beta_\ell\) 互异，则
	\[
		h_j(T)=\prod_\ell(T-\beta_\ell^{1/q_j}),
	\]
	而唯一的 \(q_j\) 次根仍互异。因此每个 \(\alpha_j\) 都在 \(K'\) 上可分。逐个添入这些元，每一步的最小多项式都整除一个无重根多项式，所以生成的有限扩张 \(LK'/K'\) 可分。

	最后核对环的有限性。每个 \(a_j^{1/q}\) 都满足次数 \(q\) 的首一方程，所以 \([k':k]\leq q^s<\infty\)。取 \(k'/k\) 的一组基，则该基乘以
	\[
		z_1^{r_1}\cdots z_d^{r_d},\qquad 0\leq r_i<q,
	\]
	生成 \(C'\) 作为 \(C\) 模，因为 \(z_i^q=x_i\)。\(z_1,\ldots,z_d\) 在 \(k'\) 上代数独立：若有非零多项式关系，取其 \(q\) 次幂，所有系数的 \(q\) 次幂都属于 \(k\)，便得到 \(x_i\) 在 \(k\) 上的非零关系，矛盾。因此 \(C'\) 是多项式环，Noether 且整闭。
\end{proof}

这里没有把整个 \(k^{1/q}\) 添入基域。即使 \(k/k^p\) 有无穷次数，上述有限列表仍只产生有限扩张 \(k'/k\)。系数列表有限来自有理函数可以写成有限项多项式的商；这也是此处与一般形式幂级数问题的关键差别。

\begin{theorem}\label{g5:thm:finite-normalization}
	设 \(k\) 为任意域，\(A\) 为有限型 \(k\) 整代数，\(L/\operatorname{Frac}A\) 为任意有限域扩张。则 \(A\) 在 \(L\) 中的整闭包为有限 \(A\) 模。特别地，\(A^\nu\) 为有限 \(A\) 模。
\end{theorem}
\begin{proof}
	由\cref{g5:prop:normalization-polynomial-reduction}，只需证明多项式环 \(C=k[x_1,\ldots,x_d]\) 在有限扩张 \(L/K\)、\(K=\operatorname{Frac}C\)，中的整闭包 \(D\) 有限于 \(C\)。

	若 \(L/K\) 可分，直接用\cref{g5:thm:separable-integral-closure-finite}。特征零下所有有限域扩张都可分：不可约多项式的导数非零且次数更低，故与原多项式互素。若特征为 \(p>0\)，用\cref{g5:lem:finite-coefficient-root-extension}构造 \(C'\)、\(K'\) 及 \(L'=LK'\)。\(C'\) 整闭且 Noether，\(L'/K'\) 有限可分，所以 \(C'\) 在 \(L'\) 中的整闭包 \(D'\) 是有限 \(C'\) 模。

	又因 \(C'\) 是有限 \(C\) 模，有限性的传递性给出 \(D'\) 为有限 \(C\) 模。每个 \(z\in D\) 满足一个 \(C\) 系数首一方程，同一方程也表明 \(z\) 整于 \(C'\)；且 \(z\in L\subseteq L'\)，故
	\[
		D\subseteq D'.
	\]
	\(D\) 是有限 \(C\) 模 \(D'\) 的子模，\(C\) Noether，所以 \(D\) 有限。再回到原归约便得结论。
\end{proof}

有限系数开根的路线见\cite[第33节，第262--263页]{Matsumura1989}。上面的引理把两项有限性分别列出：基域只加入有限多个根，变量只需要指数有界的单项式；随后才把不可分问题放到可分扩张中处理。

\subsection{约化对象的各分支与粘合}
\label{g5:subsec:0402:05}

\begin{corollary}\label{g5:cor:reduced-finite-normalization}
	设 \(k\) 为任意域，\(A\) 为约化有限型 \(k\) 代数。则 \(A^\nu\) 为有限 \(A\) 模，因而仍是有限型 \(k\) 代数、Noether 环及正规环。若 \(A\ne0\)，则
	\[
		A^\nu\cong\prod_{\mathfrak p\in\operatorname{Min}A}
		(A/\mathfrak p)^\nu.
	\]
\end{corollary}
\begin{proof}
	有限型性使 \(A\) Noether，极小素理想只有有限多个。由\cref{g5:prop:total-quotient-components}，其正规化是各整环分支的正规化的有限乘积。每个因子由\cref{g5:thm:finite-normalization}为有限 \(A/\mathfrak p\) 模，因而也为有限 \(A\) 模。有限乘积作为模等于有限直和，所以 \(A^\nu\) 有限于 \(A\)。模生成元也是代数生成元，故 \(A^\nu\) 为有限型 \(k\) 代数；Hilbert 基定理给出 Noether 性，正规性则已由\cref{g5:prop:normalization-localization}证明。零环情形直接成立。
\end{proof}

以两条轴为例，令 \(A=k[x,y]/(xy)\)。其两个分支的坐标环是 \(k[x]\)、\(k[y]\)，都整闭，所以
\[
	A^\nu=k[x]\times k[y].
\]
原环在乘积中的像是
\[
	\{(f(x),g(y)):f(0)=g(0)\}.
\]
一个多项式模 \(xy\) 可唯一写成 \(c+xf_1(x)+yg_1(y)\)，这就证明了像的描述。正规化去掉两条轴在原点的取值相等条件，把它们作为两个独立分支保存。相应的乘积环由 \(A\) 与幂等元 \((1,0)\) 生成，是一个具体的有限扩张。

\ProseParagraphBreak

现在把仿射构造放进开集。对代数闭域 \(k\) 上的约化仿射代数集 \(X\)，令
\(\widetilde X\) 为坐标环 \(k[X]^\nu\) 对应的仿射代数集。有限型与约化性保证这是第一章的映射对应所允许的坐标环，包含
\(k[X]\subseteq k[X]^\nu\) 给出有限映射
\[
	\nu:\widetilde X\longrightarrow X.
\]
在每个不可约分支的函数域上，它诱导同构：对应整闭包位于原函数域中并包含原坐标环，取分式域不改变。于是正规化保留各分支的有理函数，改变其允许作为正则函数的整元。

\ProseParagraphBreak

若 \(f\in k[X]\)，主开集 \(D(f)\) 的坐标环为 \(k[X]_f\)。局部化公式给出
\[
	(k[X]^\nu)_f=(k[X]_f)^\nu.
\]
因此先在 \(X\) 上正规化再取 \(\nu^{-1}D(f)\)，与先取 \(D(f)\) 再正规化，得到同一张仿射图。若 \(f\) 在某个分支上恒为零，这张主开图会同时删去该分支；公式已把这一情况包含在内。

\ProseParagraphBreak

这一相容性也给出一般代数簇的粘合构造。取约化代数簇的仿射开覆盖 \(X=\bigcup_i U_i\)，逐张构造 \(\widetilde U_i\)。交集 \(U_i\cap U_j\) 可用两张图的主开集作公共细化。具体地，在给定点附近先取 \(U_i\) 中的主开集 \(D(f)\subseteq U_j\)，再取 \(U_j\) 中的主开集 \(D(g)\subseteq D(f)\)。把 \(g\) 限制到 \(D(f)\) 后，可写成 \(a/f^n\)，其中 \(a\in k[U_i]\)。于是这个 \(D(g)\) 在 \(U_i\) 中恰为 \(D(fa)\)，确实同时是两张图的主开集。相应局部化后的整闭包都由同一个开集的正则函数环决定，所以它们之间有唯一的、保持原有分支函数域的同构。

\ProseParagraphBreak

在三张图的交上，这些同构满足余循环条件，因为两次复合与直接同构在每个分支函数域中都是恒等。将不交并 \(\coprod_i\widetilde U_i\) 按这些开集同构识别，取商拓扑；一个开集上的函数规定为在各仿射图上正则且在重叠处相等的函数。开集同构及余循环条件保证每张 \(\widetilde U_i\) 仍是所得空间的开图，函数的限制与粘合同样相容。各图上的 \(\nu_i\) 也在重叠处相等，因而拼成 \(\nu:\widetilde X\to X\)。更换覆盖后取公共细化，局部化公式再次给出唯一同构，所以这个构造不依赖覆盖选择。后面的概形语言将用同一组局部环与限制同态表达这一构造。

% 纲要条目排为编号小节标题，并进入目录与书签；说明段仍保留于纲要注释中。

% !TeX root = ../../Book5.tex
% 纲要标题：### 4.3 导子理想及显式计算
% 纲要位置：几何学计划纲要.md，第 5169 行。
% 正文状态：初稿；尚未正式审核。

% 纲要内容开始
% 
% * 有限双有理扩张的导子
% * 尖点 k[t²,t³] 的正规化
% * 结点的两个原像和导子
% * 商模长度、δ 不变量和局部差异
% 
% 纲要内容结束

\section{导子理想及显式计算}
\label{g5:sec:0403}

正规化把坐标环 \(A\) 放进同一个总分式环中的整闭环 \(A^\nu\)。原来只能作为分式出现的函数，可能因此成为正规函数；尖点的参数 \(t=t^3/t^2\) 就是这样的元素。要测量正规化增加了多少函数，可以考察 \(A\) 模 \(A^\nu/A\)；要确定这些新函数在哪些点附近已经属于原环，则要找出能够同时把它们乘回 \(A\) 的元素。

\ProseParagraphBreak

本节用导子理想同时回答这两个问题，再完整计算尖点和结点。计算只使用多项式、局部化和第三章的模长度。曲线正规化商的有限性可对照\cite[第4章第4.1节，定理4.11的证明，第271页]{Shafarevich2013Basic1}；用局部商的长度定义奇点不变量，以及结点、普通尖点的取值，可对照\cite[第IV章，习题1.8，第298页]{Hartshorne1977}。

\subsection{有限双有理扩张的导子}
\label{g5:subsec:0403:01}

设 \(A\) 为非零约化 Noether 环，\(Q(A)\) 为把全部非零因子变成单位所得的总分式环。若
\[
	A\subseteq B\subseteq Q(A),\qquad B\text{ 为有限 }A\text{ 模},
\]
就称这个包含为\term[youxianshuangyoulikuozhang]{有限双有理扩张}[finite birational extension]。这里的双有理性由同一个总分式环内的包含表达：若 \(S\) 为 \(A\) 的非零因子集，则 \(S^{-1}B=Q(A)\)。在整环情形，这就是两个环具有同一分式域。\cref{g5:cor:reduced-finite-normalization}保证约化有限型域代数的正规化属于这一情形。

\begin{definition}\label{g5:def:conductor}
	设 \(A\) 为非零约化 Noether 环，\(A\subseteq B\subseteq Q(A)\) 为有限双有理扩张。定义
	\[
	\mathfrak c_{B/A}=(A:B)
	=\{a\in A\mid aB\subseteq A\},
	\]
	称为这个扩张的\term[daozilixiang]{导子理想}[conductor ideal]。环 \(B\) 明确时简记为 \(\mathfrak c\)。
\end{definition}
\symidx{\(\mathfrak c_{B/A}=(A:B)\)：有限双有理扩张的导子理想}

\begin{proposition}\label{g5:prop:conductor-properties}
	设 \(A\) 为非零约化 Noether 环，\(A\subseteq B\subseteq Q(A)\) 为有限双有理扩张，\(\mathfrak c=\mathfrak c_{B/A}\)。则
	\[
	\mathfrak c=\operatorname{Ann}_A(B/A)
	=\{b\in B\mid bB\subseteq A\}.
	\]
	它同时是 \(A\) 与 \(B\) 的理想，而且是 \(B\) 中包含于 \(A\) 的最大理想。它含有一个 \(A\) 的非零因子 \(s\)，从而 \(A_s=B_s\)。此外，\(\mathfrak c=A\) 当且仅当 \(A=B\)。
\end{proposition}
\begin{proof}
	对 \(a\in A\)，条件 \(a(B/A)=0\) 就是 \(aB\subseteq A\)，得到零化子等式。若 \(b\in B\) 满足 \(bB\subseteq A\)，代入 \(1\in B\) 便有 \(b\in A\)，所以两个集合描述相同。

	对 \(c\in\mathfrak c\)、\(b\in B\)，有 \((bc)B\subseteq cB\subseteq A\)，故 \(bc\in\mathfrak c\)。加法封闭性直接成立，所以 \(\mathfrak c\) 是 \(B\) 理想，也当然是 \(A\) 理想。若 \(J\) 为 \(B\) 的理想且 \(J\subseteq A\)，则每个 \(c\in J\) 都满足 \(cB\subseteq J\subseteq A\)，因此 \(J\subseteq\mathfrak c\)。

	取有限组 \(A\) 模生成元 \(b_1,\ldots,b_r\) 生成 \(B\)。因 \(B\subseteq Q(A)\)，可写 \(b_i=a_i/s_i\)，其中 \(a_i\in A\)、\(s_i\) 为非零因子。令 \(s=\prod_i s_i\)，则 \(s\) 仍是非零因子，且 \(sb_i\in A\) 对所有 \(i\) 成立。由生成性知 \(sB\subseteq A\)，故 \(s\in\mathfrak c\)。在反演 \(s\) 后，每个 \(b\in B\) 都写为 \((sb)/s\in A_s\)，所以 \(A_s=B_s\)。最后，\(1\in\mathfrak c\) 等价于 \(B\subseteq A\)，给出最后的断言。
\end{proof}

这里 \(B/A\) 是 \(A\) 模的商；\(A\) 通常不是 \(B\) 的理想，所以这个符号并不表示商环。在可约对象上，共同分母也必须选非零因子：只有这样的元素才能同时在每个分支的函数域中反演。有限生成性使有限多个分母的乘积足以控制全部新函数。

\begin{proposition}\label{g5:prop:conductor-localization}
	设 \(A\) 为非零约化 Noether 环，\(A\subseteq B\subseteq Q(A)\) 为有限双有理扩张，\(T\subseteq A\) 为乘法集。则
	\[
	T^{-1}\mathfrak c_{B/A}
	=(T^{-1}A:T^{-1}B).
	\]
	特别地，对 \(\mathfrak p\in\operatorname{Spec}A\)，记
	\(B_{\mathfrak p}=(A\setminus\mathfrak p)^{-1}B\)，便有
	\[
	A_{\mathfrak p}=B_{\mathfrak p}
	\quad\Longleftrightarrow\quad
	\mathfrak c_{B/A}\not\subseteq\mathfrak p.
	\]
\end{proposition}
\begin{proof}
	令 \(M=B/A\)。局部化正合性给出
	\(T^{-1}M=(T^{-1}B)/(T^{-1}A)\)。一个一般的有限模 \(M\) 满足
	\[
	T^{-1}\operatorname{Ann}_A(M)
	=\operatorname{Ann}_{T^{-1}A}(T^{-1}M).
	\]
	其中从左到右的包含直接成立。为证反向，设 \(a/t\) 杀掉 \(T^{-1}M\)，并取有限组生成元 \(m_1,\ldots,m_r\) 生成 \(M\)。由局部化的零元判据，对每个 \(i\) 存在 \(u_i\in T\) 使 \(u_i a m_i=0\)。令 \(u=\prod_i u_i\)，则 \(uaM=0\)，且 \(a/t=(ua)/(ut)\) 属于左端。将这个等式用于 \(M=B/A\)，再用\cref{g5:prop:conductor-properties}便得导子的局部化公式。

	当 \(T=A\setminus\mathfrak p\) 时，两个局部化环相等，等价于它们的导子是整个 \(A_{\mathfrak p}\)，也等价于 \(\mathfrak c A_{\mathfrak p}=A_{\mathfrak p}\)。后者恰表示 \(\mathfrak c\) 含有一个不属于 \(\mathfrak p\) 的元素。
\end{proof}

这里的 \(B_{\mathfrak p}\) 只反演底环 \(A\setminus\mathfrak p\)。当 \(\mathfrak p\) 上方有多个素理想时，它可以有多个极大理想；它与在 \(B\) 的某个指定素理想处局部化是两个不同的环。

\begin{corollary}\label{g5:cor:normalization-isomorphism-locus}
	设 \(k\) 为域，\(A\) 为非零约化有限型 \(k\) 代数，\(B=A^\nu\) 为其在 \(Q(A)\) 中的正规化，\(\mathfrak c=(A:B)\)。则
	\[
	\operatorname{Supp}_A(B/A)=V(\mathfrak c),
	\qquad
	U=\{\mathfrak p\mid A_{\mathfrak p}=B_{\mathfrak p}\}
	=\operatorname{Spec}A\setminus V(\mathfrak c).
	\]
	对每个 \(f\in\mathfrak c\)，有 \(A_f=B_f\)；这些主开集 \(D_A(f)\) 覆盖 \(U\)，给出正规化在 \(U\) 上的同构。\(U\) 是具有这一性质的最大开集，并且稠密。
\end{corollary}
\begin{proof}
	由\cref{g5:cor:reduced-finite-normalization}，\(B/A\) 是有限 \(A\) 模。将\cref{g5:prop:support-annihilator}与导子的零化子描述结合，得到支集公式；上一个命题给出 \(U\) 的表达式。对 \(f\in\mathfrak c\)，包含 \(fB\subseteq A\) 直接给出 \(A_f=B_f\)。在两个这样的主开集的交 \(D_A(fg)\) 上，环同构都来自原包含的局部化，因而彼此相容。这里的同构具体指这些主开集上的坐标环相等。
	若某个点附近已经由原包含给出同构，则在该点局部化仍相等，该点必属于 \(U\)，故 \(U\) 最大。最后，\(\mathfrak c\) 含有非零因子 \(s\)。由\cref{g5:prop:total-quotient-components}，\(s\) 不属于任何极小素理想，因此 \(D_A(s)\) 在每个不可约分量中稠密；它包含于 \(U\)，故 \(U\) 稠密。
\end{proof}

\subsection{\texorpdfstring{尖点 \(k[t^2,t^3]\) 的正规化}{尖点 k[t²,t³] 的正规化}}
\label{g5:subsec:0403:02}

第一章中，尖点参数化揭示了有理逆 \(t=y/x\) 在原点不能正则延拓。现在可以把缺少的函数、阻止同构的理想以及原点上方的纤维放在同一次计算中。

\begin{example}\label{g5:exm:cusp-conductor}
	设 \(k\) 为任意域，令
	\[
	A=k[t^2,t^3]\subseteq B=k[t],\qquad x=t^2,\quad y=t^3.
	\]
	则 \(B\) 是 \(A\) 的正规化，且
	\[
	\mathfrak c_{B/A}=t^2k[t]=(t^2,t^3)_A,
	\qquad B/A\cong A/(t^2,t^3)\cong k.
	\]
\end{example}
\begin{proof}
	先核对坐标环。代入同态 \(k[x,y]\rightarrow k[t]\) 的核为 \((y^2-x^3)\)：模去这个首一的 \(y\) 二次式后，任意多项式可写成 \(F(x)+yG(x)\)，代入得到
	\[
	F(t^2)+t^3G(t^2).
	\]
	第一项只含偶次幂，第二项只含奇次幂，所以其和为零迫使 \(F=G=0\)。故
	\[
	A\cong k[x,y]/(y^2-x^3).
	\]
	又因每个整数 \(n\geq2\) 都能写成 \(2i+3j\)，有
	\begin{equation}\label{g5:eq:cusp-ring-linear-condition}
	A=k+t^2k[t].
	\end{equation}
	于是 \(B=A+At\)，并且 \(t=y/x\in\operatorname{Frac}(A)\) 满足首一方程 \(T^2-x=0\)。扩张有限双有理。由\cref{g5:prop:ufd-integrally-closed}，\(k[t]\) 整闭；若 \(z\in k(t)\) 在 \(A\) 上整，则同一首一方程也说明它在 \(B\) 上整，从而 \(z\in B\)。所以 \(B\) 正是正规化。

	由\eqref{g5:eq:cusp-ring-linear-condition}，\(t^2B\subseteq A\)，且 \(t^2B\) 为 \(B\) 理想，因此 \(t^2B\subseteq\mathfrak c\)。反过来，若 \(a=c+t^2h(t)\in\mathfrak c\)，则 \(at=ct+t^3h(t)\in A\)。\(A\) 中没有非零一次项，所以 \(c=0\)。故 \(\mathfrak c=t^2B\)。由 \(B=A+At\)，这个导子作为 \(A\) 理想由 \(t^2,t^3\) 生成。
	同一等式也表明 \(B/A\) 是由 \([t]\) 张成的一维 \(k\) 向量空间。\(t^2,t^3\) 都杀掉 \([t]\)，而常数按通常标量作用，故 \(B/A\cong A/(t^2,t^3)\) 为 \(A\) 模同构。
\end{proof}

令 \(\mathfrak m=(t^2,t^3)_A\)，它对应原点。导子就是 \(\mathfrak m\)，所以正规化只在这个点不能成为局部同构。在 \(D_A(x)\) 上，\(t=y/x\) 给出明确的逆。原点的纤维环为
\[
	B/\mathfrak mB=k[t]/(t^2).
\]
它只有素理想 \((t)/(t^2)\)，故原点在正规化中只有一个原像 \(t=0\)。这个纤维环的长度为二，而正规化商 \(B/A\) 的长度为一：前者记录纤维的厚度，后者记录原环缺少的函数方向。

\ProseParagraphBreak

把原点局部化，\(B_{\mathfrak m}\) 只有极大理想 \((t)\)，从而等于 \(k[t]_{(t)}\)。这里使用了有限整扩张的极大理想对应与刚才的纤维计算。于是原点处的正规化是一个局部环，单个参数 \(t\) 描述这一分支；原环中最低的非恒定项却从 \(t^2\) 开始。

\subsection{结点的两个原像和导子}
\label{g5:subsec:0403:03}

对结点，参数化把两个不同参数值送到同一个点。因而原环中的函数必须在这两个参数值取相同的值。这个条件可以直接写成多项式子环，进而算出正规化与导子。

\begin{example}\label{g5:exm:node-conductor}
	设 \(k\) 为特征不等于 \(2\) 的域，令
	\[
	A=k[x,y]/\bigl(y^2-x^2(x+1)\bigr),\qquad B=k[t].
	\]
	代入 \(x=t^2-1\)、\(y=t(t^2-1)\) 给出包含 \(A\subseteq B\)，并有
	\begin{equation}\label{g5:eq:node-ring-value-condition}
	A=k+(t^2-1)k[t]
	=\{f\in k[t]\mid f(1)=f(-1)\}.
	\end{equation}
	\(B\) 是 \(A\) 的正规化；原点 \(\mathfrak m=(x,y)_A\) 上方恰有 \(t=1\)、\(t=-1\) 两个点，而且
	\[
	\mathfrak c_{B/A}=(t^2-1)k[t]=(x,y)_A,
	\qquad B/A\cong A/\mathfrak m\cong k.
	\]
\end{example}
\begin{proof}
	写 \(h=t^2-1\)。先验证核理想。模去首一的 \(y\) 二次式 \(y^2-x^2(x+1)\)，每个多项式都可化为 \(F(x)+yG(x)\)。代入后为
	\[
	F(t^2-1)+t(t^2-1)G(t^2-1).
	\]
	前一项只含偶次幂，后一项只含奇次幂。若两者之和为零，则各项为零。代入 \(x\mapsto t^2-1\) 在 \(k[x]\) 上单射，因为非零多项式代入后仍有非零最高项；再用 \(th\ne0\)，可得 \(F=G=0\)。所以核恰为给定主理想，\(A\cong k[h,th]\)；特别地，\(A\) 是整环。

	为证明第一个子环等式，将任意多项式按奇偶幂分成
	\(F(t^2)+tG(t^2)\)。因 \(t^2=h+1\)，有
	\[
	hk[t]=h\,k[h]+th\,k[h]\subseteq k[h,th].
	\]
	反向包含来自 \(h,th\in hB\)，以及 \(hB\) 是 \(B\) 的理想。因此 \(A=k+hB\)。任意 \(f\in B\) 除以 \(h\) 的余式为 \(c+dt\)；条件 \(f(1)=f(-1)\) 等价于 \(2d=0\)。因特征不为 \(2\)，它等价于 \(d=0\)，即 \(f\in k+hB\)。这证明\eqref{g5:eq:node-ring-value-condition}。

	由 \(t^2=x+1\in A\)，有 \(B=A+At\)；由 \(t=y/x\)，两个环具有同一分式域。\(t\) 满足首一方程 \(T^2-(x+1)=0\)，故扩张有限且整。\(B=k[t]\) 整闭，任何在 \(A\) 上整的分式又在 \(B\) 上整，所以 \(B\) 是正规化。
	因为 \(hB\subseteq A\) 是 \(B\) 理想，有 \(hB\subseteq\mathfrak c\)。反过来，对 \(a\in\mathfrak c\)，设其在 \(1,-1\) 的共同值为 \(c\)。还须有 \(at\in A\)，即 \(c=-c\)，所以 \(c=0\)。于是 \(a\) 被 \((t-1)(t+1)=h\) 整除，故 \(\mathfrak c=hB\)。又因 \(B=A+At\)，有 \(hB=hA+thA=(x,y)_A\)。

	两个点的取值差给出满的 \(k\) 线性映射
	\[
	B\longrightarrow k,\qquad f\longmapsto f(1)-f(-1).
	\]
	满性可由 \(t/2\) 的像为 \(1\) 看出，其核由\eqref{g5:eq:node-ring-value-condition}恰为 \(A\)。把 \(k\) 作为 \(A/\mathfrak m\) 模，这个映射也是 \(A\) 线性的：若 \(a(1)=a(-1)=c\)，则 \((af)(1)-(af)(-1)=c(f(1)-f(-1))\)。因此 \(B/A\cong A/\mathfrak m\)。最后，
	\[
	B/\mathfrak mB=k[t]/(t^2-1)\cong k\times k,
	\qquad [f]\longmapsto(f(1),f(-1)),
	\]
	中国剩余定理给出这个同构。其两个素理想恰对应于 \((t-1)\)、\((t+1)\)，所以原点上方正好有这两个点。
\end{proof}

在 \(D_A(x)\) 上仍有 \(t=y/x\)，所以正规化在原点以外成为同构。原点附近，\(A\) 中的函数只记录两个参数值的共同取值；\(B\) 则允许独立指定这两个值。多出来的一个取值差就是 \(B/A\) 中的一维方向。这个结点的两点参数模型也可对照\cite[第I章第4节，例4.9.1，第29--30页]{Hartshorne1977}。

\ProseParagraphBreak

\(B_{\mathfrak m}=(A\setminus\mathfrak m)^{-1}B\) 有两个极大理想，进一步分别局部化得到 \(k[t]_{(t-1)}\) 与 \(k[t]_{(t+1)}\)。有限多个极大理想的环称为\term[banjubuhuan]{半局部环}[semilocal ring]。此处 \(B_{\mathfrak m}\) 本身仍是整环，因为它是 \(k[t]\) 的局部化；两个分支的局部环并不使它变成两个环的直积。与此相应，原来的 \(A\) 和 \(A_{\mathfrak m}\) 也都是整环，结点的两个原像不等于原局部环有两个极小素理想。

\ProseParagraphBreak

特征限制可以从方程直接核对。特征为 \(2\) 时，\(1=-1\)，并且令 \(u=t-1\) 就有 \(t^2-1=u^2\)、\(t(t^2-1)=u^3+u^2\)，所以同一个子环成为 \(k[u^2,u^3]\)。两个原像此时合成尖点模型中的一个原像，\eqref{g5:eq:node-ring-value-condition}中的两点取值条件也失去约束。

\subsection{\texorpdfstring{商模长度、\(\delta\) 不变量和局部差异}{商模长度、\GeometryPdfDelta{} 不变量和局部差异}}
\label{g5:subsec:0403:04}

尖点缺少一个一次项，结点缺少一个独立取值差；两者都表现为正规化商中的一维空间。要把这样的计算推广到任意曲线上的闭点，先证明这个商总有有限长度。

\begin{proposition}\label{g5:prop:curve-normalization-length}
	设 \(k\) 为任意域，\(A\) 为非零约化有限型 \(k\) 代数，且 \(\dim A\leq1\)。令 \(B=A^\nu\) 为正规化。则 \(M=B/A\) 是有限长度 \(A\) 模，也是有限维 \(k\) 向量空间。它的支集由有限多个极大理想组成。
\end{proposition}
\begin{proof}
	由\cref{g5:cor:reduced-finite-normalization}，\(B\) 是有限 \(A\) 模，因此 \(M\) 有限生成。\cref{g5:prop:conductor-properties}给出一个非零因子 \(s\in A\)，使 \(sM=0\)。我们先证明 \(A/(s)\) 是有限维 \(k\) 代数。

	若 \(s\) 为单位，则 \(M=0\)，结论成立。否则，每个包含 \(s\) 的素理想 \(\mathfrak p\) 都不是极小素理想，因为非零因子不属于 \(A\) 的任何极小素理想。若 \(\mathfrak p\) 还严格包含于某个素理想 \(\mathfrak q\)，在 \(\mathfrak p\) 下方选一个极小素理想，就得到长度为二的素链，与 \(\dim A\leq1\) 矛盾。故包含 \(s\) 的素理想都极大，\(A/(s)\) 为零维 Noether 环。

	由\cref{g5:thm:artinian-zero-dimension}，\(A/(s)\) 作为自身的模具有有限长度。它的组成因子都是某些剩余域 \(A/\mathfrak m\)。这些域同时是有限型 \(k\) 代数，故由\cref{g5:lem:zariski}为 \(k\) 的有限扩张。因此各组成因子的 \(k\) 维数都有限，沿组成列相加就得到 \(\dim_k A/(s)<\infty\)。
	由于 \(sM=0\)，有限生成性给出一个满同态 \((A/(s))^r\rightarrow M\)，故 \(M\) 有限维，也有有限长度。其支集包含于 \(V(s)=\operatorname{Spec}(A/(s))\)；后者由零维 Noether 环定理知为有限个闭点。
\end{proof}

在这个证明中，维数一使非零因子的零点集只剩有限个闭点。正规化增加的函数因而集中在这些点上，可以逐点计数。

\begin{definition}\label{g5:def:delta-invariant}
	设 \(k\) 为域，\(A\) 为非零约化有限型 \(k\) 代数，\(\dim A\leq1\)，\(B=A^\nu\)，并取闭点 \(\mathfrak p\in\operatorname{Spec}A\)。记
	\[
	R=A_{\mathfrak p},\qquad
	R^\nu=(A\setminus\mathfrak p)^{-1}B.
	\]
	本节采用局部长度约定，将
	\[
	\delta_{\mathfrak p}(A)
	=\ell_R(R^\nu/R)
	\]
	称为这个点的\term[deltabubianliang]{\(\delta\) 不变量}[delta invariant]。由\cref{g5:prop:normalization-localization}，正规化与局部化相容，所以 \(R^\nu\) 是 \(R\) 在其总分式环中的正规化。
\end{definition}
\symidx{\(\delta_{\mathfrak p}(A)\)：曲线点处正规化商的局部长度}

这个定义只取决于局部 \(k\) 代数 \(R\)：一个局部环同构会把非零因子、总分式环和首一方程一起对应，从而也把正规化商对应起来。有限长度模的长度为零当且仅当模为零，因此
\[
	\delta_{\mathfrak p}(A)=0
	\quad\Longleftrightarrow\quad
	A_{\mathfrak p}=B_{\mathfrak p}.
\]
原环与正规化在该点附近相同，恰好对应于 \(\delta\) 为零。

\begin{proposition}\label{g5:prop:delta-residue-degree}
	设 \(k\) 为域，\(A\) 为非零约化有限型 \(k\) 代数，\(\dim A\leq1\)，\(B=A^\nu\)，\(\mathfrak p\) 为闭点。采用\cref{g5:def:delta-invariant}中的局部长度约定，则
	\[
	\dim_k(B/A)_{\mathfrak p}
	=[\kappa(\mathfrak p):k]\,\delta_{\mathfrak p}(A),
	\]
	以及
	\[
	\dim_k(B/A)
	=\sum_{\mathfrak p\in\operatorname{Supp}_A(B/A)}
	[\kappa(\mathfrak p):k]\,\delta_{\mathfrak p}(A).
	\]
	特别地，若 \(k\) 为代数闭域，则所有剩余域次数均为一。
\end{proposition}
\begin{proof}
	局部环 \(R=A_{\mathfrak p}\) 上的简单模都是 \(\kappa(\mathfrak p)\)。对 \((B/A)_{\mathfrak p}\) 取组成列，它的每个组成因子作为 \(k\) 向量空间都具有维数 \([\kappa(\mathfrak p):k]\)。沿组成列相加，即得第一个公式。

	为证第二个公式，对有限长度 \(A\) 模 \(B/A\) 取组成列，其各因子写为 \(A/\mathfrak m_i\)。在 \(\mathfrak p\) 处局部化时，\(\mathfrak m_i=\mathfrak p\) 的因子变成 \(\kappa(\mathfrak p)\)，其余因子变成零：若 \(\mathfrak m_i\ne\mathfrak p\)，可选 \(a\in\mathfrak m_i\setminus\mathfrak p\)，它一面杀掉该因子，一面在局部化中成为单位。局部化正合性与长度加性表明，\(\delta_{\mathfrak p}\) 恰是 \(\mathfrak p\) 在这些 \(\mathfrak m_i\) 中出现的次数。把各因子的 \(k\) 维数相加并按极大理想分组，就得到第二个公式。
\end{proof}

一般域上，局部长度与对基域的向量空间维数因此需要用剩余域次数换算。\(R^\nu/R\) 本身也未必自然地是 \(\kappa(\mathfrak p)\) 向量空间；长度计数的是组成列中的剩余域因子。

\begin{proposition}\label{g5:prop:conductor-colength}
	设 \(k\) 为域，\(A\) 为非零约化有限型 \(k\) 代数，\(\dim A\leq1\)，\(B=A^\nu\)。取闭点 \(\mathfrak p\)，并令 \(R=A_{\mathfrak p}\)、\(R^\nu=B_{\mathfrak p}\)、\(\mathfrak c_R=(R:R^\nu)\)。则有短正合列
	\[
	0\longrightarrow R/\mathfrak c_R
	\longrightarrow R^\nu/\mathfrak c_R
	\longrightarrow R^\nu/R\longrightarrow0,
	\]
	其中三项都是有限长度 \(R\) 模，并且
	\[
	\delta_{\mathfrak p}(A)
	=\ell_R(R^\nu/\mathfrak c_R)
	-\ell_R(R/\mathfrak c_R).
	\]
\end{proposition}
\begin{proof}
	导子同时是两环的理想，所以包含与取商给出所示正合列。由\cref{g5:prop:conductor-localization}，\(\mathfrak c_R\) 含有全局共同分母 \(s\) 的像。上个有限性证明已经说明 \(A/(s)\) 有限长度，故 \(R/\mathfrak c_R\) 有限长度。\(R^\nu\) 是有限 \(R\) 模，所以 \(R^\nu/\mathfrak c_R\) 是有限 \(R/\mathfrak c_R\) 模，也有有限长度。最后应用\cref{g5:prop:length-additivity}。
\end{proof}

在尖点与结点中，导子都等于原点的极大理想，所以 \(R/\mathfrak c_R=k\)。正规化模去导子后分别是 \(k[t]/(t^2)\) 与 \(k\times k\)，它们作为 \(R\) 模都具有长度二。因此两者都有 \(\delta=2-1=1\)。前者保留同一点的一个一次项，后者保留两个点的独立取值，商模长度相同并不使这两个局部结构相同。

\ProseParagraphBreak

为记录正规化分离出多少个点，以下取 \(k\) 为代数闭域。对约化仿射曲线的闭点 \(\mathfrak p\)，记正规化中位于它上方的点数为 \(r_{\mathfrak p}\)，称为正规化意义下的\term[quxianfenzhishu]{分支数}[number of branches]。有限性使这个数有限，lying-over 使它至少为一。这里数的是正规化上方的点；结点模型已经表明，它可以不同于原局部环的极小素理想个数。

\begin{proposition}\label{g5:prop:delta-branch-bound}
	设 \(k\) 为代数闭域，\(A\) 为约化仿射曲线的坐标环，\(B=A^\nu\)，\(\mathfrak p\) 为闭点。则
	\[
	\delta_{\mathfrak p}(A)\geq r_{\mathfrak p}-1.
	\]
\end{proposition}
\begin{proof}
	记 \(R=A_{\mathfrak p}\)、\(R^\nu=B_{\mathfrak p}\)，并设 \(R^\nu\) 的极大理想为 \(\mathfrak n_1,\ldots,\mathfrak n_r\)。它们对应于 \(\mathfrak p\) 的全部原像，故 \(r=r_{\mathfrak p}\)。各剩余域都是 \(k\)：\(B\) 有限于 \(A\)，相应纤维是有限维 \(k\) 代数，而 \(k\) 代数闭。
	由中国剩余定理，取值映射 \(R^\nu\rightarrow\prod_iR^\nu/\mathfrak n_i=k^r\) 满射。\(R\) 的像是对角线 \(k\)，因为它在每个原像点的值都来自同一个 \(R/\mathfrak pR=k\)，且常数可以任意取。因此得到满的 \(k\) 线性映射
	\[
	R^\nu/R\longrightarrow k^r/k.
	\]
	右端维数为 \(r-1\)，左端维数由\cref{g5:prop:delta-residue-degree}等于 \(\delta_{\mathfrak p}\)，故不等式成立。
\end{proof}

\begin{table}[htbp]
	\centering
	\caption[尖点与结点的正规化数据]{尖点与结点的正规化数据；\(k\) 代数闭且特征不为 \(2\)。}
	\label{g5:tab:cusp-node-conductor}
	\begin{tabular}{lll}
	\toprule
	原点附近所比较的量 & 尖点 & 结点\\
	\midrule
	子环 \(A\subseteq k[t]\) & \(k+t^2k[t]\) & \(k+(t^2-1)k[t]\)\\
	导子 \(\mathfrak c\) & \(t^2k[t]\) & \((t^2-1)k[t]\)\\
	原点的纤维环 & \(k[t]/(t^2)\) & \(k\times k\)\\
	正规化上方点数 \(r\) & \(1\) & \(2\)\\
	正规化商长度 \(\delta\) & \(1\) & \(1\)\\
	\bottomrule
	\end{tabular}
\end{table}

\cref{g5:tab:cusp-node-conductor}把两种差异分开：结点的 \(\delta=1\) 已由分离两个点所需的一个独立取值差解释；尖点只有一个原像，\(\delta=1\) 来自该点处缺少的一次项。

\begin{example}\label{g5:exm:reducible-node-conductor}
	设 \(k\) 为任意域，取可约曲线环
	\[
	A=k[x,y]/(xy),\qquad B=k[u]\times k[v].
	\]
	令 \(x\mapsto(u,0)\)、\(y\mapsto(0,v)\)、\(1\mapsto(1,1)\)。这个代入的核恰为 \((xy)\)：模去 \((xy)\) 后的多项式写成 \(c+xF(x)+yG(y)\)，其像为 \((c+uF(u),c+vG(v))\)，为零就迫使全部系数为零。因此
	\[
	A=\{(f,g)\in k[u]\times k[v]\mid f(0)=g(0)\}.
	\]
	\cref{g5:prop:total-quotient-components}给出 \(Q(A)=k(u)\times k(v)\)。令 \(e=(1,0)\)，则 \(e^2-e=0\)，且任意 \((f,g)\) 减去 \((f(0)-g(0))e\) 后属于 \(A\)。故 \(B=A+Ae\) 为有限整扩张。两因子 \(k[u],k[v]\) 都整闭，逐分量检查首一方程便知 \(B\) 在 \(Q(A)\) 中整闭，所以它就是正规化。
	包含于 \(A\) 的 \(B\) 理想
	\[
	J=uk[u]\times vk[v]
	\]
	等于导子。确实，若 \(a=(f,g)\) 属于导子，乘以 \(e\) 后的 \((f,0)\) 仍须属于 \(A\)，所以 \(f(0)=0\)；再由 \(a\in A\) 得 \(g(0)=0\)，即 \(a\in J\)。反向由 \(J\) 的理想性质立即成立。作为 \(A\) 理想，\(J=(x,y)\)。取值差 \((f,g)\mapsto f(0)-g(0)\) 又给出 \(B/A\cong k\)，因此原点处 \(\delta=1\)，且有两个正规化原像。
	这个模型的两条全局分支本来就由两个极小素理想表示；\cref{g5:exm:node-conductor}中的曲线则整体不可约。正规化上方的点数把这两类模型在交点处共同具有的两分支现象提取了出来。
\end{example}
% 纲要条目排为编号小节标题，并进入目录与书签；说明段仍保留于纲要注释中。

% !TeX root = ../../Book5.tex
% 纲要标题：### 4.4 有限双有理映射的限制
% 纲要位置：几何学计划纲要.md，第 5176 行。
% 正文状态：初稿；尚未正式审核。

% 纲要内容开始
% 
% * 正规化与同胚、同构的区别
% * 正规化和基域扩张
% * 非约化环的约化与正规化
% * Nagata 环与正规化的有限性
% 
% ---
% 
% 纲要内容结束

\section{有限双有理映射的限制}
\label{g5:sec:0404}

尖点与结点的正规化都由直线给出，但它们在原点上方的表现不同：尖点仍只有一个原像，结点却有两个。即使点没有分开，正规化也可能增加原来不存在的正则函数。下面把点集拓扑、函数环、基域和有限性分开考察，使“有限双有理”这一条件的作用落实到可计算的模型上。

\subsection{正规化与同胚、同构的区别}
\label{g5:subsec:0404:01}

在代数闭域上的仿射代数集之间，有限映射是闭映射。因此，若某个正规化映射在点集上还是双射，它便是 Zariski 拓扑的同胚。同构还要求逆映射正则；这一要求由函数环决定。

\begin{example}[尖点的同胚正规化]
\label{g5:exm:cusp-homeomorphism}
	设 \(k\) 为代数闭域，令
	\[
		C=V(y^2-x^3)\subseteq\mathbb A_k^2,\qquad
		\nu:\mathbb A_k^1\longrightarrow C,\quad t\longmapsto(t^2,t^3).
	\]
	由\cref{g5:exm:cusp-conductor}，相应环包含
	\[
		A=k[t^2,t^3]\subseteq B=k[t]
	\]
	是有限双有理扩张，且 \(B\) 是 \(A\) 的正规化。这里也可直接看到有限性：\(t\) 满足 \(T^2-t^2=0\)，并且 \(B=A+At\)。

	映射 \(\nu\) 是点集上的双射。若 \((x,y)\in C\) 且 \(x\ne0\)，唯一可能的原像为 \(t=y/x\)，并且
	\[
		t^2=\frac{y^2}{x^2}=x,\qquad
		t^3=\frac{y^3}{x^3}=y.
	\]
	若 \(x=0\)，则 \(y=0\)，唯一原像为 \(t=0\)。由\cref{g5:prop:finite-affine-closed-fibres}，\(\nu\) 还是闭映射，故为同胚。

	它却不是同构。若逆映射正则，则直线的坐标函数 \(t\) 将来自 \(C\) 上的整体正则函数；而
	\[
		t\notin k[t^2,t^3],
	\]
	因为右侧每个多项式的一次项都为零。甚至在原点附近，\(t\) 也不能写成 \(P/Q\)，其中 \(P,Q\in A\)、\(Q(0)\ne0\)：等式 \(tQ=P\) 的左侧有非零一次项，右侧没有。这样，连续的逆公式在原点恰好缺少正则性。
\end{example}

这个例子适用于任意特征；它所需的是 \(x\ne0\) 处的除法公式，而非开平方或开立方的选择。它也是\cite[第I章，第3节，习题3.2(a)，第21页]{Hartshorne1977}中区分同胚与同构的基本模型。

\ProseParagraphBreak

对结点，取代数闭域 \(k\)、\(\operatorname{char}k\ne2\)，以及
\[
	N=V\bigl(y^2-x^2(x+1)\bigr),\qquad
	\nu(t)=\bigl(t^2-1,t(t^2-1)\bigr).
\]
由\cref{g5:exm:node-conductor}，这仍是有限双有理正规化。原点上方恰有 \(t=1,-1\) 两点；在 \(x\ne0\) 处，逆公式仍是 \(t=y/x\)。所以这次正规化会分开一个点的两条分支，底空间映射不再单射。尖点缺少一个正则函数而不缺少点；结点的函数环则把两个原像处的取值联系在一起。

\begin{proposition}
\label{g5:prop:finite-birational-normal-target}
	设 \(A\subseteq B\) 为整环，\(B\) 是有限 \(A\) 模，且包含诱导分式域相等
	\[
		\operatorname{Frac}A=\operatorname{Frac}B=K.
	\]
	若 \(A\) 在 \(K\) 中整闭，则 \(A=B\)。因此，代数闭域上两个不可约约化仿射代数集之间的有限双有理映射，若靶的坐标环整闭，就是同构。
\end{proposition}
\begin{proof}
	有限性使每个 \(b\in B\) 对 \(A\) 整。双有理性又使 \(b\in K\)。整闭性的定义于是给出 \(b\in A\)，故 \(B\subseteq A\)，与原包含合起来得到相等。几何陈述由仿射代数集与坐标环的对应得出。
\end{proof}

这一论证也解释了\cite[第2章，第5.2节，定理2.21，第129页]{Shafarevich2013Basic1}中有限双有理映射与正规化之间的关系：同一个函数域内的有限中间环必包含在整闭包中；若靶已经正规，这个空间便没有增加函数的余地。

\ProseParagraphBreak

陈述中的整环和共同分式域条件保证两侧讨论的是同一个约化不可约对象。对约化 Noether 环，可把对应条件写成
\[
	A\subseteq B\subseteq Q(A),
\]
其中 \(Q(A)\) 是总分式环；若 \(B\) 整于 \(A\)，且 \(A\) 在 \(Q(A)\) 中整闭，同样逐元素得到 \(B=A\)。这一说法保留了各分支的泛点信息。仅仅要求在靶的某个稠密开集上同构则还不够：例如
\[
	k[t]\longrightarrow k[t]\times k,\qquad
	f\longmapsto(f,f(0))
\]
有限，反演 \(t\) 后成为同构，但第二个因子给出了一个额外不可约分量。上环不是整环，两个环也没有相同的总分式环。

\subsection{正规化和基域扩张}
\label{g5:subsec:0404:02}

从域 \(k\) 扩张到域 \(K\) 时，有限型 \(k\) 代数 \(A\) 变成
\[
	A_K=A\otimes_k K.
\]
若 \(A=k[x_1,\ldots,x_n]/I\)，则 \(A_K\cong K[x_1,\ldots,x_n]/IK[x_1,\ldots,x_n]\)：方程保留，系数允许来自更大的域。因而需要检验的是这些方程在新域上的环结构，尤其是是否出现幂零元、是否仍为整环以及是否整闭。

\ProseParagraphBreak

若 \(A\subseteq B\) 是有限的 \(k\) 代数扩张，则 \(A_K\subseteq B_K\) 仍是有限扩张。单射性来自域上向量空间的张量正合性；原来的一组有限 \(A\) 模生成元张量后又生成 \(B_K\)。对应的环图为
\[
	\begin{tikzcd}
		A \arrow[r,hook] \arrow[d] & B \arrow[d] \\
		A_K \arrow[r,hook] & B_K.
	\end{tikzcd}
\]
有限性因此可以直接传到下方，而下方右环是否正规仍需另行判断。

\begin{proposition}
\label{g5:prop:normalization-field-extension}
	设 \(k\subseteq K\) 为域扩张，\(A\) 为有限型 \(k\) 整环，\(B=A^\nu\) 为其正规化。若 \(A_K\) 与 \(B_K\) 均为整环，且 \(B_K\) 在自己的分式域中整闭，则
	\[
		(A_K)^\nu=B_K
	\]
	成立，其中两侧通过自然包含放在同一个分式域中。
\end{proposition}
\begin{proof}
	由\cref{g5:thm:finite-normalization}，\(B\) 是有限 \(A\) 模。取有限组模生成元 \(b_1,\ldots,b_m\)。它们都在 \(\operatorname{Frac}A\) 中，所以有共同的非零分母 \(a\in A\) 使 \(ab_i\in A\)。从而
	\[
		B_a=A_a.
	\]
	张量 \(K\) 并与局部化交换，得到
	\[
		(B_K)_a=(A_K)_a.
	\]
	\(a\) 在 \(A_K\) 中仍非零；两个环都是整环，故反演 \(a\) 不改变分式域，于是它们有相同的分式域 \(F\)。

	又 \(B_K\) 有限于 \(A_K\)，故整于 \(A_K\)，从而 \(B_K\subseteq(A_K)^\nu\)。反过来，若 \(z\in F\) 整于 \(A_K\)，同一个首一方程也表明 \(z\) 整于 \(B_K\)。由 \(B_K\) 的整闭性，\(z\in B_K\)，得到相等。
\end{proof}

例如，对任意域扩张 \(K/k\)，尖点环的包含扩张成
\[
	K[t^2,t^3]\subseteq K[t].
\]
右侧仍是整闭整环，故尖点的正规化与这一基域扩张交换。下一例表明，一般正规 \(k\) 代数未必保留这种性质。

\begin{example}[不完美域扩张产生幂零元]
\label{g5:exm:inseparable-normality-loss}
	设 \(p\) 为素数，\(k=\mathbb F_p(s)\)，其中 \(s\) 为超越元。取
	\[
		A=k[T]/(T^p-s).
	\]
	我们先证明它是域。\(s\) 不是 \(k\) 中的 \(p\) 次幂：若 \(s=(f/g)^p\)，其中非零 \(f,g\in\mathbb F_p[s]\)，则 \(f^p=sg^p\)，比较次数得到
	\[
		p\deg f=1+p\deg g,
	\]
	不可能成立。取代数闭包中的 \(\alpha\) 满足 \(\alpha^p=s\)。在该闭包上，\(T^p-s=(T-\alpha)^p\)。如果它在 \(k[T]\) 中有次数 \(d\) 满足 \(1\leq d<p\) 的首一因子，该因子只能是 \((T-\alpha)^d\)，其 \(T^{d-1}\) 系数为 \(-d\alpha\)。因 \(d\ne0\) 于 \(k\)，这会迫使 \(\alpha\in k\)，矛盾。因此 \(T^p-s\) 不可约，\(A\) 是域，当然正规。

	现在把基域扩为 \(K=k(\alpha)\)。有
	\[
		A_K
		\cong K[T]/\bigl((T-\alpha)^p\bigr)
		\cong K[\varepsilon]/(\varepsilon^p),
		\qquad \varepsilon=T-\alpha.
	\]
	这里 \(\varepsilon\ne0\)，因为 \(1,\varepsilon,\ldots,\varepsilon^{p-1}\) 是一组 \(K\) 基；但 \(\varepsilon^p=0\)。所以 \(A_K\) 不约化，因而不是正规环。其素谱只有一个点，长度却为 \(p\)：基域扩张把原先的域变成了带有幂零厚度的点。
\end{example}

这个例子的基域不完美，所加的扩张纯不可分。一般域上，“正规”描述原来的局部环；“任意基域扩张后仍正规”是更强的要求。对有限型 \(k\) 代数，后一条件称为\term[jihezhenggui]{几何正规}[geometrically normal]：对每个域扩张 \(K/k\)，环 \(A_K\) 都正规。把“正规”换成“约化”，就得到\term[jiheyuehua]{几何约化}[geometrically reduced]。\cref{g5:exm:inseparable-normality-loss}中的域 \(A\) 连几何约化也不满足。

\subsection{非约化环的约化与正规化}
\label{g5:subsec:0404:03}

对一般环 \(A\)，每个素理想都包含幂零根式 \(\sqrt{(0)}\)，所以商映射
\[
	A\longrightarrow A_{\mathrm{red}}\coloneqq A/\sqrt{(0)}
\]
在素谱上诱导同胚。它保留点及其闭集，却会删去非零幂零元。正规环的各素局部环都是整环，因而首先必须约化；这一要求决定了下面的约定。

\begin{definition}
\label{g5:def:normalization-nonreduced}
	设 \(A\) 为 Noether 环。讨论其正规化时，先取约化环 \(A_{\mathrm{red}}\)，再取它在总分式环 \(Q(A_{\mathrm{red}})\) 中的整闭包。记作
	\[
		A^\nu\coloneqq(A_{\mathrm{red}})^\nu,
		\qquad
		A\longrightarrow A_{\mathrm{red}}\hookrightarrow A^\nu.
	\]
	因此，“先取约化”属于本书对非约化对象使用正规化时的约定；若 \(A\) 已约化，这就是\cref{g5:def:normalization-ring}中的构造。
\end{definition}

由于第二个箭头单射，复合同态的核恰为 \(\sqrt{(0)}\)。在有限型 \(k\) 代数情形，由\cref{g5:cor:reduced-finite-normalization}，\(A^\nu\) 有限于 \(A_{\mathrm{red}}\)，也就有限于 \(A\)。有限性在这里保留，但原环到正规化的映射可能失去单射性。

\begin{example}[双直线与两种整闭操作]
\label{g5:exm:double-line-reduction}
	设 \(k\) 为域，
	\[
		R=k[x,y]/(y^2).
	\]
	每个元素唯一写成 \(a(x)+y b(x)\)。若它幂零，则模 \(y\) 后的 \(a(x)\) 在整环 \(k[x]\) 中幂零，故 \(a=0\)。因此
	\[
		\sqrt{(0)}=(y),\qquad
		R_{\mathrm{red}}=k[x],\qquad
		R^\nu=k[x].
	\]
	最后一个等式来自 \(k[x]\) 整闭。复合 \(R\rightarrow R^\nu\) 把非零元 \(y\) 送到零，故不是单射。素谱上的映射却是同胚：双直线与约化直线具有相同的点和闭集。

	若改为直接在原环的总分式环中取整闭包，会得到另一种对象。\(a+yb\) 为非零因子当且仅当 \(a\ne0\)，因而
	\[
		Q(R)=k(x)[y]/(y^2).
	\]
	直接在这个环中求对 \(R\) 整的元素，所得环为
	\[
		C=k[x]+y\,k(x).
	\]
	事实上，若 \(a+yb\) 整于 \(R\)，其中 \(a,b\in k(x)\)，模 \(y\) 的首一方程说明 \(a\) 整于 \(k[x]\)，故 \(a\in k[x]\)。反过来，对任意 \(b\in k(x)\)，元素 \(yb\) 都满足 \(T^2=0\)，因而整于 \(R\)；再加上 \(a\in k[x]\) 仍为整元素。

	这个 \(C\) 既不约化，也不是有限 \(R\) 模。若有有限组生成元，可选共同分母 \(0\ne d\in k[x]\)，使每个生成元的 \(y\) 系数乘以 \(d\) 后都在 \(k[x]\) 中。于是它们的每个 \(R\) 线性组合也有这一性质。但 \(y/x^n\in C\) 对所有 \(n\geq1\) 都成立，而不可能有固定的非零多项式 \(d\) 被所有 \(x^n\) 整除。这给出矛盾。
\end{example}

双直线的两个结果分别回答不同的问题：\(R^\nu\) 从约化结构出发，得到正规环；\(C\) 则在保留原幂零元的总分式环中加入整元素。后者的无限分母现象也是\cite[第33节，第264页]{Matsumura1989}讨论非约化环整闭包时的基本机制。

\ProseParagraphBreak

这一约定还使基域扩张的反例更加明确。在\cref{g5:exm:inseparable-normality-loss}中，\(A^\nu=A\)，故
\[
	A^\nu\otimes_k K=K[\varepsilon]/(\varepsilon^p),
	\qquad
	(A\otimes_k K)^\nu=K.
\]
左侧先正规化再扩张基域，右侧先扩张基域再按约定正规化。两个环不同，差别恰好是扩张后新出现的幂零结构。

\subsection{Nagata 环与正规化的有限性}
\label{g5:subsec:0404:04}

有限型条件在正规化有限性定理中控制了分母和系数的数量。由\cref{g5:thm:finite-normalization}，任意域上的有限型整代数，在分式域的任意有限扩张内都有有限的整闭包；约化有限型代数则按有限多个分支应用这一结论。为了区分只控制原分式域与同时控制全部有限域扩张的条件，采用以下名称。

\begin{definition}
\label{g5:def:nagata-ring}
	设 \(A\) 为整环，\(F=\operatorname{Frac}A\)。
	\begin{enumerate}
		\item 若 \(A\) 在 \(F\) 中的整闭包为有限 \(A\) 模，则称 \(A\) 满足\term[N1tiaojian]{N-1 条件}[N-1 condition]。
		\item 若对每个有限域扩张 \(L/F\)，\(A\) 在 \(L\) 中的整闭包都是有限 \(A\) 模，则称 \(A\) 满足\term[N2tiaojian]{N-2 条件}[N-2 condition]，也称 \(A\) 为 Japanese 环。
	\end{enumerate}
	\foreignidx{Japanese ring}
	设 \(R\) 为交换含幺环。若 \(R\) 为 Noether 环，并且对每个素理想 \(\mathfrak p\)，整环 \(R/\mathfrak p\) 都满足 N-2 条件，则称 \(R\) 为\term[Nagatahuan]{Nagata 环}[Nagata ring]。\footnote{永田雅宜（ながた まさよし，Masayoshi Nagata，1927--2008），日本数学家，研究交换代数、代数几何与不变量理论，并构造了 Hilbert 第十四问题的反例。}
\end{definition}

N-2 蕴含 N-1，因为可取 \(L=F\)。N-1、N-2 的定义本身不要求 Noether；Nagata 条件则明确包含 Noether 性，并对全部素商同时作要求。特别地，整闭整环自动满足 N-1，因为它的整闭包就是自身；这一个事实尚未处理其他有限域扩张。Nagata 环的上述定义可对照\cite[第33节，第264页]{Matsumura1989}。

\begin{proposition}
\label{g5:prop:nagata-normalization-finite}
	任意域 \(k\) 上的有限型 \(k\) 代数都是 Nagata 环。若 \(R\) 为约化 Nagata 环，则它在 \(Q(R)\) 中的正规化是有限 \(R\) 模。
\end{proposition}
\begin{proof}
	有限型 \(k\) 代数 \(R\) 为 Noether 环。每个素商 \(R/\mathfrak p\) 仍是有限型 \(k\) 整代数，由\cref{g5:thm:finite-normalization}的有限域扩张版本满足 N-2 条件，故 \(R\) 为 Nagata 环。

	现在设 \(R\) 约化且 Nagata。零环情形直接成立，以下设 \(R\ne0\)。Noether 性使极小素理想只有有限个，记为 \(\mathfrak p_1,\ldots,\mathfrak p_r\)。由\cref{g5:prop:total-quotient-components}，
	\[
		Q(R)\cong\prod_{i=1}^r F_i,\qquad
		F_i=\operatorname{Frac}(R/\mathfrak p_i).
	\]
	令 \(C_i\) 为 \(R/\mathfrak p_i\) 在 \(F_i\) 中的整闭包。Nagata 条件给出各 \(C_i\) 有限于 \(R/\mathfrak p_i\)，因而也有限于 \(R\)。所以有限直积
	\(C=\prod_i C_i\)
	为有限 \(R\) 模，并且整于 \(R\)。

	若 \(z\in Q(R)\) 整于 \(R\)，将它的首一方程投影到第 \(i\) 个因子，便知 \(z_i\) 整于 \(R/\mathfrak p_i\)，故 \(z_i\in C_i\)。因此 \(R^\nu\subseteq C\)；反向包含来自刚才的整性，于是 \(R^\nu=C\)，正规化有限。
\end{proof}

下面给出一个完全显式的非有限正规化。它的原环不是 Noether 环，因而同时指出了该模型的适用位置。

\begin{example}[无限代数扩张造成非有限正规化]
\label{g5:exm:nonfinite-normalization}
	取无限次的代数域扩张 \(K/k\)，即每个元素都代数但 \([K:k]=\infty\)；例如可取 \(\overline{\mathbb F}_p/\mathbb F_p\)。在 \(K[x]\) 内令
	\[
		A=k+xK[x]
		=\{f(x)\in K[x]:f(0)\in k\},
		\qquad B=K[x].
	\]
	这是整环包含，且
	\[
		\operatorname{Frac}A=K(x):
	\]
	对任意 \(c\in K\)，有 \(c=(cx)/x\in\operatorname{Frac}A\)，故该分式域包含 \(K\) 和 \(x\)；反向包含由 \(A\subseteq K[x]\) 得到。

	每个 \(c\in K\) 都满足 \(k\) 系数的首一方程，因而整于 \(A\)。又 \(x\in A\) 且 \(B=A[K]\)，整元素对加法和乘法封闭，所以 \(B\) 整于 \(A\)。\(B=K[x]\) 整闭：若互素的 \(u,v\in K[x]\) 使 \(u/v\) 整于 \(K[x]\)，清去首一方程的分母可得 \(v\mid u^n\)，从而 \(v\) 是单位。因此，任何整于 \(A\) 的 \(K(x)\) 中元素也整于 \(B\)，必属于 \(B\)。这证明
	\[
		A^\nu=B=K[x].
	\]

	然而 \(B\) 不是有限 \(A\) 模。取 \(\mathfrak m=xK[x]\subset A\)，则
	\[
		A/\mathfrak m=k,\qquad
		B/\mathfrak mB=K.
	\]
	若 \(B\) 有限于 \(A\)，模去 \(\mathfrak m\) 后，有限组生成元应给出 \(K\) 在 \(k\) 上的有限张成集，与 \([K:k]=\infty\) 矛盾。

	原环 \(A\) 的 Noether 性也可以直接判定。因为
	\[
		\mathfrak m^2=x^2K[x],\qquad
		\mathfrak m/\mathfrak m^2\cong K
	\]
	作为 \(A/\mathfrak m=k\) 向量空间成立，\(\mathfrak m\) 不可能有限生成。因此 \(A\) 非 Noether。这个环的每个整元素各自只用有限个系数，但所有系数合起来需要无限个模生成元。
\end{example}

对一般 Noether 整环，正规化的有限性也确实会失败，甚至可发生在一维 Noether 局部整环上。这项反例存在性结果见\cite[第33节，第262、264页]{Matsumura1989}；它与上例的非 Noether 构造是不同的结论。因此，“Noether”“N-1”“N-2”承担不同的要求。对本章的代数几何对象，有限型定理已经提供所需有限性；推广到底环时，Nagata 条件则同时控制全部素商及其有限域扩张。

% 纲要条目排为编号小节标题，并进入目录与书签；说明段仍保留于纲要注释中。


\begin{chaptersummary}
	一维 Noether 局部整环的整闭性把极大理想压缩为一个生成元。这个统一化元使每个非零分式唯一写成单位乘整数次幂，阶随乘积相加，非零分式理想也由一个整数确定。Dedekind 整环把各点的这些整数组织成非零理想的唯一分解；可逆理想可以在每点都为主理想，却没有整体生成元。

	\ProseParagraphBreak

	约化环的总分式环同时保留各极小分支的函数域，正规化逐分支取整闭包。有限型代数通过 Noether 正规化归约到多项式基环：有限可分扩张的迹矩阵给出全部整元的共同分母；不可分扩张只需给有限多个系数及变量开根，就能进入有限可分扩张。所得整闭包有限，且构造与局部化相容，因此可以按仿射开图粘合。

	\ProseParagraphBreak

	导子是原环中能把正规化全部乘回原环的元素，也是正规化商模的零化子。它的补开集正是正规化成为同构的位置。曲线上的商模具有有限长度；局部 \(\delta\) 记录增加的组成因子数，对一般基域换成向量空间维数时还要乘剩余域次数。尖点和结点都有 \(\delta=1\)，但前者增加同一点的一次项，后者分开两个原像的取值。

	\ProseParagraphBreak

	有限双有理映射在正规靶上受整闭性约束；没有这个条件，同胚仍可能不是同构。基域扩张会改变约化性与正规性，非约化对象则先取约化再谈本章的正规化。有限型代数的有限性证明依赖明确的代数条件，不能仅凭 Noether 性推广到所有环。
\end{chaptersummary}

\BookExercises

\begin{exercise}\label{g5:ex:dvr-residue-degree-length}
	设 \(k\) 为域，\(f\in k[t]\) 为次数 \(d\geq1\) 的不可约多项式，\(R=k[t]_{(f)}\)。给定 \(0\ne q\in k[t]\)，写 \(q=f^nh\)，其中 \(f\nmid h\)。计算 \(\ell_R(R/qR)\) 和 \(\dim_k(R/qR)\)，并说明当 \(d>1\) 且 \(f\mid q\) 时两者为何不同。
\end{exercise}
\begin{solution}
	\(h\) 在 \(R\) 中为单位，所以 \(qR=f^nR\)。若 \(n=0\)，商模为零。若 \(n\geq1\)，幂滤过
	\[
		R/f^nR\supset fR/f^nR\supset\cdots
		\supset f^{n-1}R/f^nR\supset0
	\]
	有 \(n\) 个相邻商，每个都同构于剩余域
	\(\kappa=R/fR\cong k[t]/(f)\)。因此长度为 \(n\)，\(k\) 维数为 \(nd\)。也可直接识别
	\(R/f^nR\cong k[t]/(f^n)\)：与 \(f\) 互素的多项式模 \(f^n\) 已可逆，局部化不再改变商环。长度把一个 \(\kappa\) 因子计为一，而它作为 \(k\) 向量空间有 \(d\) 维。
\end{solution}

\begin{exercise}\label{g5:ex:finite-dvr-degree}
	设 \(R\subseteq S\) 为 DVR 之间的有限扩张，且结构同态为局部同态。记分式域为 \(K\subseteq L\)，统一化元为 \(\pi,\tau\)，剩余域为 \(\kappa\subseteq\lambda\)，并写
	\[
		\pi=u\tau^e,\qquad u\in S^\times,\quad e\geq1.
	\]
	先证明 DVR 上每个有限生成无挠模都自由，再证明
	\[
		[L:K]=e[\lambda:\kappa].
	\]
\end{exercise}
\begin{solution}
	设 \(M\) 为有限生成无挠 \(R\) 模。取 \(M/\pi M\) 的一组基并提升为 \(m_1,\ldots,m_r\)，由 Nakayama 引理，它们生成 \(M\)。若有非平凡关系 \(\sum_i a_im_i=0\)，从全部非零系数中取最小阶 \(n\)，写 \(a_i=\pi^nb_i\)，至少一个 \(b_i\) 为单位。无挠性允许约去 \(\pi^n\)，再模去 \(\pi M\)，得到所选剩余基之间的非平凡关系，矛盾。因此这些生成元是一组自由基。

	\(S\) 为整环，故作为 \(R\) 模无挠，因而自由，设秩为 \(r\)。把 \(R\setminus\{0\}\) 反演后得到有限维 \(K\) 整代数，所以是域，恰为 \(L\)。于是 \(r=[L:K]\)，且
	\[
		\dim_\kappa(S/\pi S)=r.
	\]
	另一方面，\(\pi S=\tau^eS\)，商环的幂滤过有 \(e\) 个组成因子，每个同构于 \(\lambda\)。\(\lambda\) 是有限维 \(\kappa\) 向量空间，因为它是 \(S/\pi S\) 的商。因此
	\[
		r=\dim_\kappa(S/\tau^eS)=e[\lambda:\kappa].
	\]
	这里次数分成了统一化元的阶与剩余域扩张次数两项。
\end{solution}

\begin{exercise}\label{g5:ex:dedekind-six-factorizations}
	沿用\cref{g5:exm:dedekind-nonprincipal}中的
	\(A=\mathbb Z[\alpha]\)、\(\alpha^2=-5\)，令
	\[
		P=(2,1+\alpha),\qquad
		Q_+=(3,1+\alpha),\qquad Q_-=(3,1-\alpha).
	\]
	证明
	\[
		(2)=P^2,\quad (3)=Q_+Q_-,\quad
		(1+\alpha)=PQ_+,\quad(1-\alpha)=PQ_-.
	\]
	据此比较 \(6=2\cdot3=(1+\alpha)(1-\alpha)\) 的元素分解与理想分解。
\end{exercise}
\begin{solution}
	\(P^2=(2)\) 已在正文直接计算。模去 \(3\)，有
	\[
		A/(3)\cong\mathbb F_3[T]/(T^2-1)
		\cong\mathbb F_3\times\mathbb F_3.
	\]
	因此 \((3)=Q_+\cap Q_-=Q_+Q_-\)，两理想的商环均为 \(\mathbb F_3\)。

	\(P+Q_+=A\)，因为和中含 \(2,3\)。元素 \(1+\alpha\) 属于两者，故
	\((1+\alpha)\subseteq P\cap Q_+=PQ_+\)。中国剩余定理说明 \(A/PQ_+\) 有六个元素；而令 \(\alpha=-1\) 得
	\[
		A/(1+\alpha)\cong\mathbb Z/6\mathbb Z.
	\]
	包含关系两端的商具有相同有限大小，所以两理想相等。取共轭并注意 \((2,1-\alpha)=P\)，便得 \((1-\alpha)=PQ_-\)。

	范数 \(N(a+b\alpha)=a^2+5b^2\) 乘法相容，范数为一的元素为单位，且没有范数为二或三的元素。由范数分别为 \(4,9,6,6\)，可知 \(2,3,1+\alpha,1-\alpha\) 都不可约：任何非单位分解都会要求出现范数二或三的因子。两种元素分解的因子范数不同，因而不是调整单位或次序得到的同一分解。它们对应的理想分解却都为
	\[
		(6)=P^2Q_+Q_-,
	\]
	与 Dedekind 整环的唯一理想分解一致。
\end{solution}

\begin{exercise}\label{g5:ex:quadratic-squarefree-normal}
	设 \(k\) 为特征不为二的域，\(f(x)\in k[x]\) 为非恒定平方自由多项式，即其不可约因子的指数全为一。证明
	\[
		B=k[x,y]/(y^2-f(x))
	\]
	为整闭整环。若 \(a\in k\) 为 \(f\) 的根，证明点 \((a,0)\) 处的局部环为 DVR，且
	\[
		\operatorname{ord}(y)=1,\qquad
		\operatorname{ord}(x-a)=2.
	\]
\end{exercise}
\begin{solution}
	令 \(C=k[x]\)、\(K=k(x)\)。平方自由性使 \(f\) 不是 \(K\) 中的平方，因此 \(T^2-f\) 在 \(K[T]\) 中不可约。首一除法说明 \(B=C\oplus Cy\) 单射到域 \(L=K(y)\)，故为整环。

	任取 \(z=u+vy\in L\) 整于 \(B\)，其中 \(u,v\in K\)。因 \(B\) 有限于 \(C\)，\(z\) 也整于 \(C\)。共轭 \(u-vy\) 同样整，于是
	\[
		2u,\qquad u^2-v^2f
	\]
	都整于 \(C\)，且属于 \(K\)。\(C\) 整闭，\(2\) 又可逆，故 \(u\in C\)、\(v^2f\in C\)。若 \(v=a/b\) 为互素多项式之比，而不可约多项式 \(q\) 整除 \(b\)，则 \(b^2\mid a^2f\) 迫使 \(q^2\mid f\)，与平方自由性矛盾。因此 \(b\) 为单位，\(v\in C\)，从而 \(z\in B\)。

	若 \(f(a)=0\)，平方自由性保证 \(f=(x-a)h(x)\) 且 \(h(a)\ne0\)。在 \((a,0)\) 处，\(h\) 为单位，
	\[
		x-a=y^2h^{-1}.
	\]
	局部极大理想 \((x-a,y)\) 因而由 \(y\) 生成。它是非域 Noether 局部整环，由 DVR 刻画得 \(y\) 为统一化元；再由上式得到 \(x-a\) 的阶为二。
\end{solution}

\begin{exercise}\label{g5:ex:trace-dual-larger}
	设 \(k\) 为特征不为二的域，\(C=k[x]\subset D=k[t]\)，其中 \(x=t^2\)。令 \(K=k(x)\)、\(L=k(t)\)，定义
	\[
		D^\ast=\{z\in L:\operatorname{Tr}_{L/K}(zD)\subseteq C\}.
	\]
	求基 \(1,t\) 的迹矩阵及 \(D^\ast\)，并证明 \(D\subsetneq D^\ast\)。说明为什么迹的整值条件给出整闭包的上界，却不单独刻画整元。
\end{exercise}
\begin{solution}
	扩张为可分二次扩张，共轭把 \(t\) 送到 \(-t\)。因此
	\[
		\bigl(\operatorname{Tr}(e_ie_j)\bigr)_{i,j}
		=\begin{pmatrix}2&0\\0&2x\end{pmatrix},
		\qquad \Delta=4x,
	\]
	其中 \(e_1=1,e_2=t\)。对 \(z=a+bt\)、\(a,b\in K\)，因 \(D\) 的 \(C\) 基是 \(1,t\)，只需检查
	\[
		\operatorname{Tr}(z)=2a\in C,\qquad
		\operatorname{Tr}(zt)=2xb\in C.
	\]
	故
	\[
		D^\ast=C\oplus C\,\frac{t}{x}=C\oplus C\,\frac1t.
	\]
	元素 \(1/t\) 属于右侧却不属于 \(k[t]\)，所以包含严格。

	环 \(D=k[t]\) 整闭且整于 \(C\)，分式域为 \(L\)，所以它正是 \(C\) 在 \(L\) 中的整闭包。因此 \(1/t\) 虽满足全部迹整值条件，却不是整元。迹配对把整元限制在一个有限模中；有限性证明只需要这个包含，并没有把它当作等式。
\end{solution}

\begin{exercise}\label{g5:ex:infinite-p-rank-finite-roots}
	设 \(p\) 为素数，\(k=\mathbb F_p(s_1,s_2,\ldots)\)，其中 \(s_i\) 代数独立。令
	\[
		C=k[x],\quad K=k(x),\quad
		c^p=s_1,\quad z^p=x,\quad L=K(c,z).
	\]
	求 \(C\) 在 \(L\) 中的整闭包及其 \(C\) 模秩。说明这个计算为什么不需要给所有 \(s_i\) 同时开 \(p\) 次根。
\end{exercise}
\begin{solution}
	在 \(k\) 中，\(s_1\) 不是 \(p\) 次幂。可把 \(k\) 看成 \(\mathbb F_p(s_2,s_3,\ldots)(s_1)\)，在不可约多项式 \(s_1\) 处的阶为一，而任何 \(p\) 次幂的阶都被 \(p\) 整除。因此 \([k(c):k]=p\)。同样，\(x\) 在 \(k(c)(x)\) 中不是 \(p\) 次幂，所以 \([L:K]=p^2\)。

	令 \(D=k(c)[z]\)。这是多项式整环，因而整闭，分式域为 \(L\)。它整于 \(C\)，并且
	\[
		c^a z^b,\qquad 0\leq a,b<p,
	\]
	是一组自由 \(C\) 基。张成性由 \(c^p=s_1\)、\(z^p=x\) 得到；独立性则先按 \(z\) 的指数模 \(p\) 分组，再用 \(1,c,\ldots,c^{p-1}\) 在 \(k\) 上的独立性。故 \(D\) 是秩 \(p^2\) 的有限自由 \(C\) 模。

	若 \(u\in L\) 整于 \(C\)，它也整于 \(D\)，由 \(D\) 整闭得 \(u\in D\)。因此 \(D\) 正是所求整闭包。构造只涉及关系中出现的系数 \(s_1\)；其他 \(s_i\) 仍留在原系数域中。即使系数域有无限多个彼此独立的不可分方向，这个指定有限扩张所需的开根仍为有限次。
\end{solution}

\begin{exercise}\label{g5:ex:coordinate-axes-normalization}
	设 \(k\) 为域，\(r\geq2\)，
	\[
		A=k[x_1,\ldots,x_r]/(x_ix_j:i\ne j).
	\]
	求 \(Q(A)\)、\(A^\nu\)、导子以及原点处的 \(\delta\)。比较原点的局部环与其正规化中的极大理想个数。
\end{exercise}
\begin{solution}
	每个元素唯一写成
	\[
		c+\sum_{i=1}^r x_if_i(x_i).
	\]
	把它送到 \((c+t_if_i(t_i))_{i=1}^r\)，便把 \(A\) 单射到
	\(B=\prod_i k[t_i]\)，像恰为各分量常数项相同的元组。极小素理想是分别保留一条轴的核，其交为零，故 \(A\) 约化，且
	\[
		Q(A)=\prod_i k(t_i),\qquad A^\nu=B.
	\]
	第二式由各分支多项式环整闭得到。

	令 \(J=\prod_i t_i k[t_i]\)。这是包含于 \(A\) 的 \(B\) 理想，故包含于导子。反过来，若 \(b=(b_i)\) 属于导子，乘以第 \(i\) 个坐标幂等元后仍须具有相同常数项。因 \(r\geq2\)，这迫使 \(b_i(0)=0\)。于是导子恰为 \(J\)，作为 \(A\) 理想就是原点的极大理想 \(\mathfrak m=(x_1,\ldots,x_r)\)。

	取常数项给出
	\[
		B/A\cong k^r/k(1,\ldots,1).
	\]
	这个商被 \(\mathfrak m\) 杀掉，长度与 \(k\) 维数都为 \(r-1\)，故 \(\delta_{\mathfrak m}=r-1\)。原点局部环只有一个极大理想，正规化的半局部环 \(B_{\mathfrak m}\) 则有 \(r\) 个；其纤维环为 \(k^r\)，分别记录各条轴的原点。
\end{solution}

\begin{exercise}\label{g5:ex:semigroup-three-four}
	设 \(k\) 为域，\(A=k[t^3,t^4]\subset k[t]\)。求正规化、导子及唯一非正规点处的 \(\delta\)，并给出正规化商的一组 \(k\) 基和一条组成列。比较正规化的纤维长度与 \(\delta\)。
\end{exercise}
\begin{solution}
	\(t=t^4/t^3\) 属于 \(\operatorname{Frac}A\)，并满足 \(T^3-t^3=0\)。因此 \(B=k[t]\) 有限整于 \(A\)，而 \(B\) 整闭且与 \(A\) 有相同分式域，所以 \(A^\nu=B\)。

	指数半群 \(\langle3,4\rangle\) 为
	\[
		0,3,4,6,7,8,9,\ldots;
	\]
	从六开始的每个整数都属于它，因为先有 \(6,7,8\)，再反复加三。缺少的非负指数恰为 \(1,2,5\)，故 \(B/A\) 的 \(k\) 基为
	\[
		[t],\ [t^2],\ [t^5].
	\]
	导子为 \(t^6B\)。一向来自所有次数至少六的项都属于 \(A\)。若一个非零导子元素的最低次数 \(n<6\)，则 \(n\) 必为 \(0,3,4\) 之一；分别乘以 \(t,t^2,t\)，会产生不可抵消的最低缺项 \(t,t^5,t^5\)，与仍属于 \(A\) 矛盾。作为 \(A\) 理想，导子为 \((t^6,t^7,t^8)\)。

	导子的根式为 \(\mathfrak m=(t^3,t^4)\)，故只有这个点非正规。商模有组成列
	\[
		0\subset k\cdot[t^5]\subset k\cdot[t^2]+k\cdot[t^5]
		\subset B/A,
	\]
	这里方括号表示 \(B/A\) 中的剩余类。各子空间都被 \(A\) 的作用保持，每个相邻商均为 \(A/\mathfrak m\cong k\)。因此 \(\delta_{\mathfrak m}=3\)。

	纤维环则为
	\[
		B/\mathfrak mB=k[t]/(t^3),
	\]
	长度也为三。在本例中两个数恰好相等；尖点 \(k[t^2,t^3]\) 的纤维长度为二而 \(\delta=1\)，已经表明两者没有普遍相等关系。
\end{solution}

\begin{exercise}\label{g5:ex:identify-many-points}
	设 \(k\) 为域，\(a_1,\ldots,a_r\in k\) 两两不同，\(r\geq2\)。令
	\[
		g(t)=\prod_{i=1}^r(t-a_i),\qquad
		A=k+g(t)k[t]\subseteq B=k[t].
	\]
	证明 \(A\) 为有限型 \(k\) 整代数，\(B=A^\nu\)，并求导子、\(\delta\) 与 \(\mathfrak m=gB\) 上方的正规化纤维。说明这个原环虽为整环，仍可在同一个点上有 \(r\) 个正规化原像。
\end{exercise}
\begin{solution}
	写 \(g(t)=t^r+c_{r-1}t^{r-1}+\cdots+c_0\)。对 \(n\geq r\)，有
	\[
		gt^n=g(gt^{n-r})-\sum_{j=0}^{r-1}c_jgt^{n-r+j}.
	\]
	归纳说明所有 \(gt^n\) 都由 \(g,gt,\ldots,gt^{r-1}\) 代数生成。因此
	\[
		A=k[g,gt,\ldots,gt^{r-1}]
	\]
	为有限型整代数。又 \(t=(gt)/g\)，故分式域为 \(k(t)\)；\(t\) 满足首一方程 \(g(T)-g(t)=0\)，所以 \(B=A[t]\) 有限于 \(A\)。\(B\) 整闭，因而为正规化。

	取值同态与中国剩余定理给出
	\[
		B/gB\cong k^r.
	\]
	\(A\) 的像是对角线 \(k\)，所以
	\[
		A=\{f\in k[t]:f(a_1)=\cdots=f(a_r)\}.
	\]
	导子包含 \(gB\)。若其中一个元素模 \(gB\) 是非零对角常数 \(c\)，乘以分别在各点取任意值的插值多项式，就会使结果不再是对角常数，矛盾。因此导子恰为 \(gB=\mathfrak m\)。

	再有 \(B/A\cong k^r/k\)，被 \(\mathfrak m\) 杀掉，所以唯一非正规点的 \(\delta=r-1\)。纤维环 \(B/\mathfrak mB=k^r\)，对应于 \(t=a_1,\ldots,a_r\)。原环嵌入整环 \(k[t]\)，故自身以及原点局部化均为整环；这些原像数记录的是正规化分离的点，并非原局部环的极小素理想数。
\end{solution}

\begin{exercise}\label{g5:ex:normalization-quotient-and-lifts}
	设 \(k\) 为特征不为二的域。用 \(P=k[x,y]\)、\(I=(y^2-x^3)\) 证明正规化一般不与取商相容。再令
	\[
		A=k[t^2-1,t(t^2-1)]\subseteq B=k[t],
	\]
	证明结点原点的取值同态 \(A\to k\) 至少有两个不同的延伸 \(B\to k\)。这与正文对单射 \(A\to C\)、\(C\) 整闭时的唯一延伸有何区别？
\end{exercise}
\begin{solution}
	\(P\) 是多项式整环，已整闭，所以 \(P^\nu=P\)。但
	\[
		P/I\cong k[t^2,t^3]
	\]
	不整闭，其正规化为 \(k[t]\)。因此
	\[
		P^\nu/IP^\nu=P/I\ne(P/I)^\nu.
	\]
	取商改变了函数域和整性问题的工作环境，不能把原环的正规性直接传到商环。

	对结点环，原点同态把 \(t^2-1\) 和 \(t(t^2-1)\) 都送到零。\(B\) 上在 \(t=1\) 与 \(t=-1\) 处取值，都延伸这个同态；由于特征不为二，它们在 \(t\) 上的取值不同。这里 \(A\to k\) 有非零核，不会诱导分式域的嵌入。正文中的唯一性则通过单射把每个整元的分式表示固定在 \(\operatorname{Frac}C\) 中，因而不适用于这个取值同态。
\end{solution}

\begin{exercise}\label{g5:ex:geometrically-nonnormal-reduced}
	设 \(p\) 为奇素数，\(k=\mathbb F_p(s)\)，并令
	\[
		A=k[X,Y]/(X^2-(Y^p-s)).
	\]
	证明 \(A\) 为整闭整环。取 \(K=k(a)\)、\(a^p=s\)，证明 \(A_K=A\otimes_kK\) 仍为整环，却不整闭；求其正规化，并计算 \(X=Y-a=0\) 处的 \(\delta\)。
\end{exercise}
\begin{solution}
	由\cref{g5:exm:inseparable-normality-loss}的不可约性计算，\(Y^p-s\) 在 \(k[Y]\) 中不可约。因此 \(Y^p-s\) 虽导数为零，作为 \(k[Y]\) 中的不可约因子却只出现一次，仍是平方自由多项式。应用\cref{g5:ex:quadratic-squarefree-normal}的证明，得到 \(A\) 为整闭整环。

	扩域并令 \(Z=Y-a\)，得到
	\[
		A_K\cong K[X,Z]/(X^2-Z^p).
	\]
	代入 \(Z=t^2,X=t^p\)，可把这个环识别为 \(K[t^2,t^p]\)。确实，模去 \(X^2-Z^p\) 后，每个元素唯一写成 \(b(Z)+Xc(Z)\)；代入后的两项分别只有偶次幂与奇次幂，因为 \(p\) 为奇数，故其和为零迫使 \(b=c=0\)。因此代入单射，环仍为整环。

	在分式域中有
	\[
		t=\frac{X}{Z^{(p-1)/2}},
	\]
	它满足首一方程 \(T^2-Z=0\)，却不属于 \(K[t^2,t^p]\)，因为后者没有一次项。因此这个扩域环不整闭，正规化为 \(K[t]\)。指数半群 \(\langle2,p\rangle\) 缺少的非负整数恰为
	\[
		1,3,\ldots,p-2.
	\]
	正规化商因此有 \(K\) 维数 \((p-1)/2\)。反演 \(Z\) 后，上面的分式已给出 \(t\)，所以正规化在 \(D(Z)\) 上成为同构；而 \(V(Z)\) 中因 \(X^2=0\) 只剩原点。因此商的支集仅为原点，剩余域为 \(K\)，该点的 \(\delta=(p-1)/2\)。这个例子中，基域扩张保持约化性，却改变了正规性。
\end{solution}

\begin{exercise}\label{g5:ex:nilpotent-finite-generic-isomorphism}
	设 \(k\) 为域，\(A=k[t]\)，
	\[
		B=k[t,\varepsilon]/(t\varepsilon,\varepsilon^2).
	\]
	证明自然同态 \(A\to B\) 有限且单射，反演 \(t\) 后为同构，素谱之间也给出同胚，但原环同态不是同构。说明为什么正规靶上的有限双有理同构结论不能只用“在一个稠密开集上同构”来替代其结构假设。
\end{exercise}
\begin{solution}
	每个 \(B\) 元素唯一写为 \(f(t)+c\varepsilon\)，其中 \(f\in k[t]\)、\(c\in k\)。因此 \(A\) 单射到 \(B\)，且 \(1,\varepsilon\) 作为 \(A\) 模生成 \(B\)。反演 \(t\) 后，关系 \(t\varepsilon=0\) 迫使 \(\varepsilon=0\)，故 \(B_t=A_t\)。开集 \(D(t)\) 在 \(\operatorname{Spec}A\) 中稠密。

	每个素理想都包含幂零元 \(\varepsilon\)，且 \(B/(\varepsilon)=A\)，所以素谱按商环对应与 \(\operatorname{Spec}A\) 同胚。可是 \(\varepsilon\ne0\)，故 \(B\) 不是约化环，也不可能等于 \(A\)。特别地，\(B\) 不能嵌入 \(Q(A)=k(t)\)：域中没有非零幂零元。正文同构结论要求上环确实处在相同的分式环内，这个条件会排除只支撑在特殊点上的额外幂零信息。
\end{solution}

\begin{exercise}\label{g5:ex:real-nonsplit-normalization-fiber}
	令 \(B=\mathbb R[t]\)、\(g=t^2+1\)，并取
	\[
		A=\mathbb R+gB,\qquad \mathfrak m=gB.
	\]
	证明 \(A\) 为有限型整代数、\(B=A^\nu\)，计算导子、\(\delta_{\mathfrak m}\) 及纤维环。说明为什么这个实有理点没有实有理的正规化原像，却不违背 lying-over。
\end{exercise}
\begin{solution}
	由 \(t^2=g-1\)，反复降次得到 \(A=\mathbb R[g,gt]\)。元素 \(t\) 满足 \(T^2-(g-1)=0\)，且 \(t=(gt)/g\)，所以 \(B\) 为同一分式域中的有限扩张。\(B\) 整闭，故为正规化。

	模去 \(gB\)，有
	\[
		B/gB\cong\mathbb C,
	\]
	而 \(A\) 的像是实常数。因此 \(B/A\cong\mathbb C/\mathbb R\)，是被 \(\mathfrak m\) 杀掉的一维实向量空间。导子包含 \(gB\)；若导子元素的像为非零实数 \(c\)，乘以 \(t\) 后的像为 \(ci\notin\mathbb R\)，与乘积仍在 \(A\) 矛盾。所以导子恰为 \(\mathfrak m\)，且 \(\delta_{\mathfrak m}=1\)。

	纤维环为
	\[
		B/\mathfrak mB=\mathbb C.
	\]
	它只有一个素点，剩余域为 \(\mathbb C\)，不存在到 \(\mathbb R\) 的实代数同态，所以没有实有理点。lying-over 保证的是一个素理想收缩到 \(\mathfrak m\)，这里的 \((t^2+1)\) 正是这样的素理想；它并不要求该点的剩余域等于基域。
\end{solution}

\begin{exercise}\label{g5:ex:intermediate-normalization-conductors}
	设 \(k\) 为域，
	\[
		A=k[t^4,t^5]\subseteq B=k[t^2,t^5]\subseteq C=k[t].
	\]
	证明这两步扩张均有限，求三个导子 \((A:B)\)、\((B:C)\)、\((A:C)\)，并计算
	\(\dim_k(B/A)\)、\(\dim_k(C/B)\)、\(\dim_k(C/A)\)。
\end{exercise}
\begin{solution}
	\(B=A[t^2]\)，其中 \(t^2\) 满足 \(T^2-t^4=0\)；又 \(C=B[t]\)，其中 \(t\) 满足 \(T^2-t^2=0\)，所以两步均有限。三个环的分式域都是 \(k(t)\)，而 \(C\) 整闭，因此 \(C\) 同时是前两个环的正规化。

	半群 \(\langle4,5\rangle\) 缺少的非负指数为
	\[
		1,2,3,6,7,11;
	\]
	从十二开始的每个指数都出现，因为 \(12,13,14,15\) 出现，随后可反复加四。半群 \(\langle2,5\rangle\) 缺少的指数为 \(1,3\)，从四开始全部出现。因此
	\[
		(A:C)=t^{12}C,\qquad (B:C)=t^4C.
	\]
	这些导子等式也可逐项判断：一个次数属于原半群的单项式进入导子，当且仅当加上任意非负整数后仍属于原半群；不同单项式相乘后的次数互不混合，故可分别检查系数。

	判断 \((A:B)\) 时只须要求加上 \(\langle2,5\rangle\) 中的指数后仍属于 \(\langle4,5\rangle\)。满足条件的指数恰为
	\[
		8,10,12,13,14,15,\ldots,
	\]
	即 \(8+\langle2,5\rangle\)。所以
	\[
		(A:B)=t^8B=(t^8,t^{10})_A.
	\]
	例如指数九不满足条件，因为 \(9+2=11\) 是缺项；指数八和十则满足条件，十二以后也都满足。

	最后，各商的基由相应半群之差给出：
	\[
		\dim_k(B/A)=4,\qquad
		\dim_k(C/B)=2,\qquad
		\dim_k(C/A)=6.
	\]
	其中 \(B/A\) 可取 \([t^2],[t^6],[t^7],[t^{11}]\) 为基。短正合列
	\[
		0\longrightarrow B/A\longrightarrow C/A
		\longrightarrow C/B\longrightarrow0
	\]
	解释了 \(6=4+2\)。本例还满足
	\((A:B)(B:C)=(A:C)\)；一般只能直接保证包含于右侧，因为若 \(aB\subseteq A\)、\(bC\subseteq B\)，则 \(abC\subseteq A\)。
\end{solution}
